% Main file. Compile with pdfLaTeX (three passes), latexmk -pdf, or Tectonic.
% All illustrations are original teaching figures from tools/reproduce_examples.py.
\documentclass[11pt,a4paper,oneside,openany]{book}
\usepackage[T1]{fontenc}
\usepackage{lmodern}
\usepackage[margin=24mm,headheight=15pt]{geometry}
\usepackage{amsmath,amssymb,mathtools,bm}
\usepackage{microtype}
\usepackage{xcolor,graphicx}
\usepackage{booktabs,longtable,tabularx,array}
\usepackage{enumitem}
\usepackage{tikz}
\usetikzlibrary{arrows.meta,positioning,calc}
\usepackage{fancyhdr}
\usepackage{hyperref}
\usepackage{bookmark}
\definecolor{ink}{HTML}{17354C}
\definecolor{teal}{HTML}{087E8B}
\definecolor{rust}{HTML}{A84429}
\hypersetup{hypertexnames=false,colorlinks=true,linkcolor=ink,citecolor=teal,urlcolor=teal,
 pdftitle={Understanding the Aberration Budget Paper: A First-Principles Textbook},
 pdfauthor={A self-study companion for a beginning graduate student},
 pdfsubject={EUV lithography, wavefront aberrations, Zernike polynomials, partial coherence, and inverse design}}
\graphicspath{{figures/}}
\pagestyle{fancy}
\fancyhf{}
\fancyhead[L]{\small\color{ink}Understanding the aberration budget paper}
\fancyhead[R]{\small\thepage}
\renewcommand{\headrulewidth}{0.3pt}
\fancypagestyle{plain}{\fancyhf{}\fancyfoot[C]{\thepage}\renewcommand{\headrulewidth}{0pt}}
\setlength{\parindent}{0pt}
\setlength{\parskip}{5pt plus 1pt}
\setlength{\emergencystretch}{2em}
\setlist{itemsep=3pt,topsep=5pt}
\setcounter{tocdepth}{1}
\setcounter{secnumdepth}{2}
\allowdisplaybreaks[2]
\newcommand{\dd}{\mathrm{d}}
\newcommand{\ii}{\mathrm{i}}
\newcommand{\ee}{\mathrm{e}}
\newcommand{\vct}[1]{\bm{#1}}
\newcommand{\avg}[1]{\left\langle #1\right\rangle}
\newcommand{\abs}[1]{\left|#1\right|}
\newcommand{\norm}[1]{\left\lVert#1\right\rVert}
\newcommand{\nm}{\,\mathrm{nm}}
\newcommand{\um}{\,\mu\mathrm{m}}
\newcommand{\NA}{\mathrm{NA}}
\newcommand{\OPD}{\mathrm{OPD}}
\newcommand{\sinc}{\operatorname{sinc}}
\newcommand{\Real}{\operatorname{Re}}
\newcommand{\Imag}{\operatorname{Im}}
\newcommand{\diag}{\operatorname{diag}}
\newcommand{\Var}{\operatorname{Var}}
\newcommand{\Cov}{\operatorname{Cov}}
\newcommand{\E}{\mathbb{E}}
\newcommand{\R}{\mathbb{R}}
\newcommand{\FT}{\mathcal{F}}
\newcommand{\paperloc}[1]{\textit{Paper location: #1}}
\newcommand{\reading}[1]{\begin{quote}\small\color{teal}\textbf{Textbook connection.} #1\end{quote}}
\newcommand{\physicalmeaning}[1]{\begin{quote}\textbf{\color{teal}Physical meaning.} #1\end{quote}}
\newcommand{\careful}[1]{\begin{quote}\textbf{\color{rust}Read carefully.} #1\end{quote}}
\newcommand{\checkpoint}[1]{\begin{quote}\textbf{Check your understanding.} #1\end{quote}}
\newcommand{\worked}[1]{\subsection*{Worked example: #1}}
\newcolumntype{Y}{>{\raggedright\arraybackslash}X}
\newcolumntype{L}[1]{>{\raggedright\arraybackslash}p{#1}}
\newcommand{\figteach}[3]{\begin{figure}[tbp]\centering\includegraphics[width=#2\textwidth]{#1}\caption{#3}\end{figure}}
\begin{document}
\frontmatter
\begin{titlepage}
\color{ink}
{\large A first-principles graduate self-study textbook\par}
\vspace{18mm}
{\Huge\bfseries Understanding the\\[3mm] Aberration Budget\\[3mm] Paper\par}
\vspace{10mm}
{\Large From waves and diffraction to EUV lithography,\\[2mm]
Zernike modes, neural networks, and inverse design\par}
\vspace{12mm}
\rule{\textwidth}{0.8pt}
\vspace{5mm}
{\large A detailed companion to\par}
\vspace{3mm}
{\large Zhishu Chen, Lisong Dong, Huwen Ding, and Yayi Wei,\par}
\textit{Aberration budget analysis of EUV lithography from the imaging
performance of a contact layer in a 5 nm technology node},\\
\textit{Applied Optics} \textbf{62}, 7270--7279 (2023).\\[3mm]
\href{https://doi.org/10.1364/AO.496461}{doi:10.1364/AO.496461}
\vfill
{\large Built around the supplied Born--Wolf and Saleh--Teich textbooks\par}
\vspace{5mm}
{\small Original explanations, worked derivations, scientific diagrams,
solved exercises, a source audit, and reproducible teaching calculations.\\
Prepared 24 September 2026.\par}
\end{titlepage}
\chapter*{How to study this book}
\addcontentsline{toc}{chapter}{How to study this book}
This book assumes no previous optics, lithography, aberration theory, or
machine-learning knowledge. Basic algebra is useful. Calculus, complex
numbers, Fourier analysis, and matrix ideas are introduced as they become
necessary. The aim is to make every technical step in Chen and colleagues'
paper \cite{chen} understandable, including the steps compressed into
phrases such as rigorous simulation, Zernike coefficient, process variation
band, and inverse sensitivity analysis.

The research paper contains only three numbered equations. Those equations
are not the whole theory behind the work. The missing chain is
\[
\text{field}\ \longrightarrow\ \text{mask diffraction}\
\longrightarrow\ \text{pupil phase}\ \longrightarrow\ \text{image}\
\longrightarrow\ \text{edges}\ \longrightarrow\ \text{constraints}.
\]
The neural network is an approximate way to travel backward along part of
this chain. Learning the forward physics first makes the backward problem
much easier to understand.

The explanations are original teaching material. Source equations are
rederived in explicit conventions rather than copied without context.
Numerical demonstrations labelled \emph{teaching examples} are calculations
made for this book. They are not recovered data from the authors' commercial
simulations. In particular, no trained version of their neural network,
mask-simulator settings file, or proprietary dataset has been supplied.

\section*{A practical study sequence}
\begin{enumerate}
\item Read Chapters \ref{ch:story}--\ref{ch:waves} with a notebook. You
should be able to convert nanometres of optical path into radians of phase
and explain why field amplitudes interfere.
\item Study Chapters \ref{ch:diffraction}--\ref{ch:zernike}. Draw a pupil,
an image plane, a defocused wavefront, and a Zernike mode. Do the small
integrals yourself.
\item Work through Chapters \ref{ch:partial}--\ref{ch:sensitivity}.
These connect optics to the three measured quantities in the paper.
\item Read the data and learning chapters, then open the paper alongside
Chapter \ref{ch:walkthrough}. Every numbered equation, table, and figure
has a corresponding explanation.
\item Finish the budget and laboratory chapters and the solved exercises.
Try each exercise before looking at its answer.
\end{enumerate}
Do not try to memorize 25 polynomial labels. Understand the pupil function,
the unit convention, and the map from phase to edges. Those ideas remain
useful when the numbering convention or lithographic pattern changes.

\section*{What was actually available in the source files}
\begin{itemize}
\item The paper is complete: 10 PDF pages, journal pages 7270--7279.
Paper p.~3 means the third PDF page, journal page 7272.
\item The supplied \emph{Principles of Optics} scan is the fourth edition
of Born and Wolf \cite{bw}, with 859 PDF pages. Its printed pagination
and PDF pagination differ, and a central run is reversed. For example,
printed pages 464, 465, and 466 are PDF pages 555, 554, and 553.
Use the verified entry points in Appendix \ref{app:reading}, then follow
the printed page numbers.
\item The file named \texttt{Saleh \& Teich Foundations of Photonics.pdf}
actually contains \emph{Fundamentals of Photonics}, second edition
\cite{st}. It has 300 PDF pages and ends at printed page 278.
The useful chapters on ray optics, waves, Fourier optics,
electromagnetic optics, and polarization are present. Chapter 11,
\emph{Statistical Optics}, is listed in the contents but is not included.
This book therefore uses Born--Wolf for the missing partial-coherence
foundation and supplies its own complete derivations.
\end{itemize}

\section*{How to interpret statements in this book}
\textbf{Reported result} means something explicitly stated or displayed
in the paper. \textbf{Derived result} follows from equations and assumptions
written here. \textbf{Teaching example} uses a stated simplified model.
\textbf{Unresolved detail} means the publication does not supply enough
information to settle it. These distinctions matter because an attractive
graph is not a complete simulation specification.

Several publication issues deserve particular attention: the pattern tables
disagree, the output count and stated Zernike range disagree, a Table 4
column header is repeated, and the source-to-pupil conventions are incomplete.
The guide preserves the reported information and identifies the uncertainty.
It does not quietly choose a correction and present it as an author's fact.

This textbook concerns the supplied 2023 study. Its conclusions are read
in the context of its patterns, illumination, simulation conditions, and
validation data. It supplies the underlying education without treating a
simulation study as a complete production qualification.

\tableofcontents
\mainmatter
\part{The physical language of the paper}
\chapter{What problem is the paper trying to solve?}\label{ch:story}
\paperloc{Title, abstract, Introduction, and Conclusion.}

\section{Printing a pattern with light}
A semiconductor chip is made by repeatedly creating patterns in thin
layers on a wafer. In optical lithography, light illuminates a patterned
mask, an optical system projects a smaller image onto a light-sensitive
material called photoresist, and chemical processing turns the exposure
into a material pattern. Etching or deposition then transfers or uses that
pattern. The paper concerns the optical exposure part of this sequence.

A \emph{contact hole} is an opening intended to connect device regions or
conducting layers. A \emph{via} is also an interconnection opening, commonly
between conducting levels. For the optical argument, both are small
two-dimensional features whose position and size must be controlled.
A square or rectangle in the mask does not automatically print as an
identical square or rectangle: diffraction rounds corners and mixes light
from nearby features.

The \emph{critical dimension}, abbreviated CD, is a specified feature width.
A rectangle needs at least a horizontal width $\mathrm{CD}_x$ and a vertical
width $\mathrm{CD}_y$. \emph{Pitch} is the centre-to-centre repetition
distance of a periodic pattern. It is not the width of a hole and not the
space between neighbouring edges. For a 14 nm square hole on 38 nm pitch,
the nominal edge-to-edge gap along an axis is $38-14=24\nm$.

The name \emph{5 nm technology node} does not mean that every printed line,
hole, or transistor dimension is 5 nm. It labels a technology generation.
The authors use nominal hole widths of 14 nm and 28 nm. The 2021 IRDS
lithography roadmap explicitly separates node labels from physical
dimensions \cite{irds}. Always identify the layer and feature before
assigning a physical length to a node name.

\section{What makes an image imperfect?}
Even a perfectly manufactured finite optical system has diffraction.
That unavoidable spreading is different from \emph{aberration}, the
departure of an actual optical wavefront from a chosen ideal wavefront.
Aberrations change the phases with which light from different pupil
locations contributes to an image.

The important objects form a hierarchy:
\begin{enumerate}
\item \textbf{Physical hardware:} mirror shapes, positions, coatings,
the mask stack, and illumination optics.
\item \textbf{Optical response:} amplitude, phase, polarization, and the
diffraction orders admitted by the pupil.
\item \textbf{Image:} a continuous distribution of optical intensity.
\item \textbf{Printed geometry:} edges selected by the resist process.
\item \textbf{Performance indicators:} widths, shifts, and changes under
focus and exposure variations.
\end{enumerate}
A neural network operating on the final indicators does not directly
specify a set of manufacturable mirror surfaces. It proposes a wavefront
description that still needs an optical realization.

\section{Why a budget is useful}
Suppose a contact must remain within $0.4\nm$ of its intended vertical
position. Many optical errors contribute to that position. A budget asks
how much of each error can be tolerated while satisfying the performance
requirement. A conventional budget might restrict the total RMS wavefront
error. This paper instead explores combinations of wavefront coefficients
that are acceptable for a particular contact-layer test set.

Write a vector of aberration coefficients as
$\vct c=(c_1,\ldots,c_K)^T$ and the computed imaging indicators as
$\vct y=F(\vct c)$. The symbol $F$ stands for an entire optical and
image-measurement calculation. The forward task calculates $\vct y$
from known $\vct c$. The inverse task finds some $\vct c$ giving acceptable
$\vct y$.

The paper trains an approximate inverse map $G$:
\[
\vct c_{\rm candidate}=G(\vct y_{\rm desired}),\qquad
F(\vct c_{\rm candidate})\approx\vct y_{\rm desired}.
\]
It then sends the proposed coefficients back through the commercial
forward simulator. That final comparison is essential: a network output
is a candidate, and its actual simulated image determines whether it works.

\physicalmeaning{The central proposal is to choose allowable wavefronts using
the imaging needs of a particular layer. The contribution is the
layer-specific inverse modelling and candidate-screening workflow,
not a new law of diffraction.}

\section{Vocabulary you will repeatedly meet}
\begin{longtable}{@{}p{.23\textwidth}p{.71\textwidth}@{}}
\toprule Term & Meaning in this study\\ \midrule\endhead
EUV & Extreme ultraviolet light; here the wavelength is 13.5 nm.\\
DUV & Deep ultraviolet; the Introduction compares EUV with 193 nm light.\\
Projector / objective & The optical system forming the mask image on
the wafer; in EUV it uses reflective optics.\\
Pupil & A representation of the aperture selecting ray directions or
spatial-frequency components.\\
Wavefront & A surface of equal optical phase at a specified time.\\
OPD / WFE & Optical path difference / wavefront error relative to a
specified ideal reference, usually expressed as a length.\\
Zernike coefficient & The amplitude of one specified polynomial mode
in a pupil-wavefront expansion.\\
PS & Pattern shift: displacement of a specified printed feature centre.\\
PVB & Process variation band: the spread of printed contours when
specified process conditions change.\\
NILS & Normalized image log slope: relative intensity steepness at an edge.\\
BPNN & Backpropagation neural network, a trainable nonlinear function.\\
OPC & Optical proximity correction: modification of the mask layout to
improve the eventual printed pattern.\\
\bottomrule
\end{longtable}
\checkpoint{If two patterns have the same CD, must they have the same
pitch, centre position, or process robustness? No. These are separate
properties; this is why the paper uses several indicators for many patterns.}

\chapter{The mathematics we actually need}\label{ch:math}
\section{Units and dimensionless ratios}
One nanometre is $10^{-9}$ metres; one micrometre is $10^{-6}$ metres.
An angle in radians is the ratio of arc length to radius and is dimensionless.
An exponential must have a dimensionless argument.
For example, $2\pi W/\lambda$ is legitimate because $W$ and $\lambda$
are both lengths.

A \emph{wave} of optical path error means one wavelength of OPD.
A milliwave, written here as $\mathrm{m}\lambda$, is $10^{-3}$ waves:
\begin{equation}
W_{\rm nm}=\lambda_{\rm nm}\frac{c_{\mathrm{m}\lambda}}{1000}.
\label{eq:milliwave}
\end{equation}
The paper's MMW is used in this milliwave sense. It is not a millimetre.
At 13.5 nm, $1\,\mathrm{m}\lambda=0.0135\nm$ and
$20\,\mathrm{m}\lambda=0.270\nm$.

\section{Complex numbers are a language for phase}
The imaginary unit obeys $\ii^2=-1$. A complex number $z=a+\ii b$ has
conjugate $z^*=a-\ii b$, modulus $|z|=\sqrt{a^2+b^2}$, and squared modulus
$zz^*=a^2+b^2$. Euler's formula is
\begin{equation}
\ee^{\ii\phi}=\cos\phi+\ii\sin\phi.
\end{equation}
Multiplication by $\ee^{\ii\phi}$ rotates a complex number by $\phi$.
It changes phase without changing magnitude.

Represent a real harmonic electric field by
\begin{equation}
E(\vct r,t)=\Real\{U(\vct r)\ee^{-\ii\omega t}\},\qquad
U=A\ee^{\ii\phi}.
\label{eq:phasor}
\end{equation}
The observable field is real. The complex quantity $U$ stores its
amplitude and phase compactly. Time differentiation becomes multiplication
by $-\ii\omega$, a major simplification.

For two scalar fields $U_1$ and $U_2$, expand the intensity:
\begin{align}
|U_1+U_2|^2
&=(U_1+U_2)(U_1^*+U_2^*)\\
&=|U_1|^2+|U_2|^2+U_1U_2^*+U_1^*U_2\\
&=A_1^2+A_2^2+2A_1A_2\cos(\phi_1-\phi_2).
\label{eq:interference}
\end{align}
The last term is interference. Equal amplitudes in phase give intensity
$4A^2$; equal amplitudes out of phase by $\pi$ give zero. The separate
intensities were $A^2$ in both cases. This is why changing phase can
change an image without changing the amount of light admitted by the pupil.

\section{Derivatives, gradients, and Taylor expansion}
A derivative measures a local change. For small displacement $\delta x$,
\[
I(x+\delta x)=I(x)+I'(x)\delta x+O(\delta x^2).
\]
Here $O(\delta x^2)$ denotes terms whose leading size is proportional
to the square of the small displacement. A partial derivative varies
one argument while holding the others fixed:
\[
\nabla_\perp W=\left(\frac{\partial W}{\partial x},
\frac{\partial W}{\partial y}\right).
\]
For $W=A(x^2+y^2)^2$,
$\nabla_\perp W=4A(x^2+y^2)(x,y)$.
If $W$ is a length, this gradient is dimensionless; it will be related
to an angular error.

We frequently use
\begin{align}
\sqrt{1+\epsilon}&=1+\frac{\epsilon}{2}
-\frac{\epsilon^2}{8}+O(\epsilon^3),\\
\ee^{\ii\phi}&=1+\ii\phi-\frac{\phi^2}{2}+O(\phi^3),\\
\ln(1+\epsilon)&=\epsilon+O(\epsilon^2).
\end{align}
These follow by evaluating successive derivatives at zero. They are
approximations for small dimensionless arguments. Calling a formula
paraxial or small-aberration specifies the expansion that has been made;
it does not remove the need to check the argument's size.

\worked{Changing to normalized pupil coordinates}
Let $x_p=a u$, $y_p=a v$, where $a$ is the physical pupil radius,
and $\widetilde W(u,v)=W(au,av)$. The chain rule gives
\begin{equation}
\frac{\partial W}{\partial x_p}
=\frac{1}{a}\frac{\partial\widetilde W}{\partial u}.
\label{eq:chain}
\end{equation}
For $\widetilde W=B u^2$, the physical slope is $2Bu/a$,
not $2Bu$. Omitting $1/a$ changes both dimensions and physics.

\section{Integrals over a disk}
Use normalized polar coordinates $u=\rho\cos\theta$,
$v=\rho\sin\theta$, $0\leq\rho\leq1$.
The small area is $\dd u\,\dd v=\rho\,\dd\rho\,\dd\theta$.
The uniformly weighted disk average is
\begin{equation}
\avg{f}=\frac{1}{\pi}\int_0^{2\pi}\int_0^1
f(\rho,\theta)\rho\,\dd\rho\,\dd\theta.
\label{eq:diskaverage}
\end{equation}
The factor $\rho$ accounts for the growing circumference of rings.
Sampling equally spaced radii with equal weights is not uniform area sampling.
For $q>-2$,
\begin{equation}
\avg{\rho^q}=2\int_0^1\rho^{q+1}\dd\rho=\frac{2}{q+2}.
\label{eq:diskpowers}
\end{equation}
Thus $\avg{\rho^2}=1/2$, $\avg{\rho^4}=1/3$, and
$\avg{\rho^8}=1/5$. Also $\avg{\rho^2\cos^2\theta}=1/4$
because the angular average of $\cos^2\theta$ is $1/2$.

The mean-subtracted RMS of a real wavefront is
\begin{equation}
\sigma_W=\sqrt{\avg{(W-\avg W)^2}}
=\sqrt{\avg{W^2}-\avg W^2}.
\label{eq:variance}
\end{equation}
It is a length. The mean $\avg W$ is called piston.
For $W=B\rho^2$, $\sigma_W=|B|/\sqrt{12}$.

\section{Vectors, matrices, and least squares}
A vector is an ordered list. A matrix applies weighted sums to that list.
If $A$ has $M$ rows and $K$ columns, $A\vct c$ has $M$ entries:
\[
(A\vct c)_i=\sum_{j=1}^K A_{ij}c_j.
\]
The transpose $A^T$ interchanges rows and columns. The Euclidean norm
obeys $\norm{\vct c}^2=\vct c^T\vct c=\sum_jc_j^2$.
To fit data $\vct b$ with $A\vct c$, minimize
$L=\tfrac12\norm{A\vct c-\vct b}^2$.
Expansion and differentiation give
$\nabla_{\vct c}L=A^T(A\vct c-\vct b)$.
At a stationary point,
\begin{equation}
A^TA\,\vct c=A^T\vct b.
\label{eq:normal}
\end{equation}
These are the normal equations. A unique solution requires independent
columns. If a nonzero $\vct v$ satisfies $A\vct v=0$, then $\vct c$
and $\vct c+\vct v$ predict the same data. Such a vector is in the
\emph{null space}. This idea becomes central to the paper's claim that
different aberration combinations can have similar imaging performance.

\section{Averages, random errors, and correlations}
For samples $x_1,\ldots,x_N$, the mean is $\bar x=N^{-1}\sum_i x_i$.
An expectation $\E[x]$ is the corresponding probability-weighted average.
Variance is $\E[(x-\E[x])^2]$ and
\[
\Cov(x,y)=\E[(x-\E[x])(y-\E[y])].
\]
Expanding the square proves
\begin{equation}
\Var(x+y)=\Var(x)+\Var(y)+2\Cov(x,y).
\label{eq:covsum}
\end{equation}
The root-sum-square rule follows only when covariance vanishes.
Two fabrication errors that always move together do not behave like
independent errors. Fixed errors do not become random simply because
their signs are unknown.

\section{Fourier analysis: a shape as a sum of ripples}
A periodic function of period $p$ can be represented as
\begin{equation}
t(x)=\sum_{m=-\infty}^{\infty}a_m\ee^{\ii2\pi mx/p},\qquad
a_m=\frac1p\int_{-p/2}^{p/2}t(x)\ee^{-\ii2\pi mx/p}\dd x.
\label{eq:fourierseries}
\end{equation}
Substitute the series into the integral. The integral of
$\ee^{\ii2\pi(m-n)x/p}$ over one period is zero unless $m=n$, when it
equals $p$. The integral selects one ripple.

For nonperiodic two-dimensional fields we use
\begin{align}
\widetilde U(\vct f)&=\int_{\R^2}U(\vct x)
\ee^{-\ii2\pi\vct f\cdot\vct x}\dd^2x,\\
U(\vct x)&=\int_{\R^2}\widetilde U(\vct f)
\ee^{\ii2\pi\vct f\cdot\vct x}\dd^2f.
\label{eq:fourierpair}
\end{align}
Spatial frequency $\vct f=(f_x,f_y)$ has units of inverse length.
Fine detail needs higher spatial frequencies. A finite pupil removes
some frequencies, so an image cannot reproduce arbitrary fine detail.

Convolution, $(h*t)(x)=\int h(x-x')t(x')\dd x'$, means summing shifted
responses to all object points. Substitution of the Fourier pair shows
$\FT[h*t]=\widetilde h\,\widetilde t$: convolution in position becomes
multiplication in frequency. This is the algebraic foundation of Fourier
optics, not a separate empirical imaging rule.

\chapter{Light, phase, rays, and optical path}\label{ch:waves}
\reading{Saleh--Teich \S\S2.1--2.2, printed pp.~40--47 (PDF pp.~62--69),
and \S2.5, p.~58 (PDF p.~80). Born--Wolf \S1.3, p.~14 (PDF p.~43),
and \S3.1, p.~109 (PDF p.~138).}

\section{From Maxwell's equations to a wave}
In a homogeneous, source-free, linear, isotropic, nonmagnetic medium,
Maxwell's equations include
\[
\nabla\times\vct E=-\mu_0\frac{\partial\vct H}{\partial t},\quad
\nabla\times\vct H=\epsilon\frac{\partial\vct E}{\partial t},\quad
\nabla\cdot\vct E=0.
\]
Take the curl of the first equation and substitute the second.
Using $\nabla\times(\nabla\times\vct E)
=\nabla(\nabla\cdot\vct E)-\nabla^2\vct E$ gives
\begin{equation}
\nabla^2\vct E-\mu_0\epsilon\frac{\partial^2\vct E}{\partial t^2}=0.
\end{equation}
The wave speed is $v=1/\sqrt{\mu_0\epsilon}=c/n$.
For one frequency and a scalar component under a justified approximation,
\begin{equation}
\nabla^2 U+n^2k_0^2U=0,\qquad
k_0=\frac{2\pi}{\lambda_0}=\frac{\omega}{c}.
\label{eq:helmholtz}
\end{equation}
This is the Helmholtz equation. A solution
$U=A\ee^{\ii\vct k\cdot\vct r}$ is a plane wave if
$|\vct k|=nk_0$. Constant-phase surfaces are planes perpendicular to
$\vct k$. given that $|\vct k|=nk_0$, $\lambda_{medium}  = \frac{2\pi}{|\vct k|} = \frac{\lambda}{n}$

\subsection*{From the Poynting vector to optical intensity}
We now derive the intensity formula rather than introduce it as a
separate optical rule. The result for a travelling plane wave in a
lossless, nonmagnetic dielectric is
\[
I=\frac{n\epsilon_0c}{2}|\vct E_0|^2.
\]
Here $\vct E_0$ is an electric-field phasor, measured in volts per metre,
and its components are defined using peak rather than RMS amplitudes.
The factor $1/2$ comes from averaging the real oscillating fields.
The factor $n\epsilon_0c$ comes from the relation between the electric
and magnetic fields. We will obtain these factors separately.

\subsubsection*{1. The Poynting vector measures energy flow}
For real, instantaneous fields in SI units, the Poynting vector is
\begin{equation}
\vct S(\vct r,t)=\vct E(\vct r,t)\times\vct H(\vct r,t)
=\frac{1}{\mu_0}\vct E(\vct r,t)\times\vct B(\vct r,t),
\label{eq:poynting-instant}
\end{equation}
where $\vct B=\mu_0\vct H$ in the nonmagnetic medium considered here.
The cross product gives a direction of energy flow. Its units are
\[
[\vct S]=\frac{\mathrm V}{\mathrm m}
          \frac{\mathrm A}{\mathrm m}
=\frac{\mathrm W}{\mathrm m^2},
\]
so it is power per area. For a surface element $\dd A$ with unit normal
$\widehat{\vct n}_A$, the signed power crossing it is
$\vct S\cdot\widehat{\vct n}_A\,\dd A$.

Why does this particular cross product describe energy?
For a lossless, nondispersive dielectric with constant $\epsilon$,
use the vector identity
\[
\nabla\cdot(\vct E\times\vct H)
=\vct H\cdot(\nabla\times\vct E)
-\vct E\cdot(\nabla\times\vct H).
\]
Substituting the source-free Maxwell equations above gives
\begin{align}
\nabla\cdot\vct S
&=-\mu_0\vct H\cdot\frac{\partial\vct H}{\partial t}
  -\epsilon\vct E\cdot\frac{\partial\vct E}{\partial t}
\notag\\
&=-\frac{\partial}{\partial t}
  \left(\frac{\epsilon|\vct E|^2}{2}
             +\frac{\mu_0|\vct H|^2}{2}\right).
\end{align}
The expression in parentheses is the electromagnetic energy density
$u_{\rm em}$, in joules per cubic metre. Thus
\begin{equation}
\frac{\partial u_{\rm em}}{\partial t}+\nabla\cdot\vct S=0,
\qquad
\frac{\dd}{\dd t}\int_Vu_{\rm em}\,\dd V
=-\oint_{\partial V}\vct S\cdot\widehat{\vct n}_A\,\dd A.
\label{eq:poynting-conservation}
\end{equation}
The second equality follows from the divergence theorem: outward energy
flow through a boundary reduces the energy stored inside it.
This identifies $\vct S$ as an energy-flux density.
The simple instantaneous formula for $u_{\rm em}$ assumes no material
dispersion; the monochromatic flux calculation below only needs the
real, lossless material parameters at the frequency of interest.

\subsubsection*{2. Average the real fields, not just their phasors}
An optical detector or resist exposure averages over many optical cycles.
Let $T_{\rm opt}=2\pi/\omega$ and define a one-cycle average at a fixed point:
\[
\avg{\vct S}_t
=\frac{1}{T_{\rm opt}}\int_{t_0}^{t_0+T_{\rm opt}}\vct S(t)\,\dd t.
\]
For an exactly monochromatic field this equals an average over any
integer number of cycles. Write the real fields in terms of their
complex spatial phasors:
\begin{align*}
\vct E(\vct r,t)
&=\Real\{\widetilde{\vct E}(\vct r)\ee^{-\ii\omega t}\}
=\frac12\left(
\widetilde{\vct E}\ee^{-\ii\omega t}
+\widetilde{\vct E}^{\,*}\ee^{+\ii\omega t}\right),\\
\vct H(\vct r,t)
&=\Real\{\widetilde{\vct H}(\vct r)\ee^{-\ii\omega t}\}
=\frac12\left(
\widetilde{\vct H}\ee^{-\ii\omega t}
+\widetilde{\vct H}^{\,*}\ee^{+\ii\omega t}\right).
\end{align*}
Suppressing the common position argument, their cross product is
\begin{align*}
\vct E\times\vct H
=\frac14\Bigl[&
(\widetilde{\vct E}\times\widetilde{\vct H})\ee^{-2\ii\omega t}
+\widetilde{\vct E}\times\widetilde{\vct H}^{\,*}\\
&+\widetilde{\vct E}^{\,*}\times\widetilde{\vct H}
+(\widetilde{\vct E}^{\,*}\times\widetilde{\vct H}^{\,*})
\ee^{+2\ii\omega t}\Bigr].
\end{align*}
The first and last terms average to zero because
$\avg{\ee^{\pm2\ii\omega t}}_t=0$.
The two remaining terms are complex conjugates of one another.
Since $z+z^*=2\Real z$, applied component by component,
\begin{equation}
\boxed{\avg{\vct S}_t
=\frac12\Real\!\left(
\widetilde{\vct E}\times\widetilde{\vct H}^{\,*}\right).}
\label{eq:poynting-average}
\end{equation}
The conjugation and the real part are essential: an unmodified cross
product of complex phasors is not the measured average power flow.
This expression also applies to monochromatic fields that are not plane
waves; the following reduction to electric-field magnitude alone uses
the plane-wave assumptions.

\subsubsection*{3. Maxwell's equations relate the two field amplitudes}
For a travelling plane wave, write
\[
\widetilde{\vct E}(\vct r)
=\vct E_0\ee^{\ii\vct k\cdot\vct r},\qquad
\widetilde{\vct H}(\vct r)
=\vct H_0\ee^{\ii\vct k\cdot\vct r},
\qquad
\vct k=k\widehat{\vct s},
\]
where $\widehat{\vct s}$ is a real unit vector in the propagation direction
and $k=n\omega/c$.
For the time convention $\ee^{-\ii\omega t}$, Faraday's law becomes
\[
\nabla\times\widetilde{\vct E}
=\ii\omega\mu_0\widetilde{\vct H}.
\]
Taking the curl of the plane-wave exponential gives
$\nabla\times\widetilde{\vct E}
=\ii\vct k\times\widetilde{\vct E}$.
Cancel the common exponential and $\ii$ to obtain
\begin{equation}
\vct H_0
=\frac{\vct k\times\vct E_0}{\omega\mu_0}
=\frac{1}{\eta}\widehat{\vct s}\times\vct E_0,
\qquad
\eta=\sqrt{\frac{\mu_0}{\epsilon}}
=\frac{\eta_0}{n},\qquad
\eta_0=\sqrt{\frac{\mu_0}{\epsilon_0}}.
\label{eq:plane-wave-impedance}
\end{equation}
We used $\epsilon=n^2\epsilon_0$ and
$c=1/\sqrt{\mu_0\epsilon_0}$.
The quantity $\eta$ is the wave impedance: it sets the ratio of electric
to magnetic field amplitude and has units of ohms.
Also, $\nabla\cdot\widetilde{\vct E}=0$ gives
$\widehat{\vct s}\cdot\vct E_0=0$, so the electric field is transverse.
The magnetic field is transverse as well, and its direction follows from
the cross product in Eq.~\eqref{eq:plane-wave-impedance}.

\subsubsection*{4. Obtain the intensity and its direction}
Because $\vct k$ is real, the spatial factors
$\ee^{\ii\vct k\cdot\vct r}$ and
$\ee^{-\ii\vct k\cdot\vct r}$ cancel in
Eq.~\eqref{eq:poynting-average}. Because $\eta$ is real,
\begin{align*}
\avg{\vct S}_t
&=\frac{1}{2\eta}
\Real\!\left[\vct E_0\times
             (\widehat{\vct s}\times\vct E_0^*)\right]\\
&=\frac{1}{2\eta}
\Real\!\left[
\widehat{\vct s}(\vct E_0\cdot\vct E_0^*)
-\vct E_0^*(\vct E_0\cdot\widehat{\vct s})\right]\\
&=\frac{|\vct E_0|^2}{2\eta}\widehat{\vct s}.
\end{align*}
The second line is the vector triple-product identity.
Transversality removes its second term, while
$|\vct E_0|^2=\vct E_0\cdot\vct E_0^*
=|E_{0x}|^2+|E_{0y}|^2+|E_{0z}|^2$ is real and nonnegative.
Define the plane-wave intensity as the average power per area
perpendicular to propagation:
\begin{equation}
\boxed{
I=\avg{\vct S}_t\cdot\widehat{\vct s}
=\frac{|\vct E_0|^2}{2\eta}
=\frac{n\epsilon_0c}{2}|\vct E_0|^2.}
\label{eq:plane-wave-intensity}
\end{equation}
For the last equality,
$1/\eta=n/\eta_0=n\sqrt{\epsilon_0/\mu_0}=n\epsilon_0c$.
The result holds for linear, elliptical, and circular polarization in
this medium; it uses the sum of the squared component phasor magnitudes.

If a detector surface normal makes angle $\alpha$ with
$\widehat{\vct s}$, its incident power per detector area is
$\avg{\vct S}_t\cdot\widehat{\vct n}_A=I\cos\alpha$, taking the normal
toward the direction of energy flow. Thus one must specify the area
orientation when converting flux into detected power.

\subsubsection*{5. Check the factor of two with a real cosine wave}
Take a linearly polarized wave travelling along $+z$:
\[
\vct E=\widehat{\vct x}E_{\rm pk}\cos(kz-\omega t),\qquad
\vct H=\widehat{\vct y}\frac{E_{\rm pk}}{\eta}
\cos(kz-\omega t).
\]
Then
\[
\vct S=\widehat{\vct z}\frac{E_{\rm pk}^2}{\eta}
\cos^2(kz-\omega t),\qquad
\avg{\cos^2(kz-\omega t)}_t=\frac12,
\]
which gives $I=E_{\rm pk}^2/(2\eta)$ directly.
If the field is instead specified by
$E_{\rm rms}=E_{\rm pk}/\sqrt2$, the same physical result is
$I=E_{\rm rms}^2/\eta$.
The two formulas use different amplitude conventions, so the factor
$1/2$ must not be inserted twice.
As a units-and-scale check, a vacuum plane wave with
$E_{\rm pk}=1\,\mathrm{V/m}$ has
$I=1/(2\eta_0)\simeq1.33\times10^{-3}\,\mathrm{W/m^2}$,
using $\eta_0\simeq376.73\,\Omega$.

\subsubsection*{6. Why scalar optics can write \texorpdfstring{$I=|U|^2$}{I = |U| squared}}
For a fixed unit polarization vector $\widehat{\vct e}$, write
$\widetilde{\vct E}=\widehat{\vct e}\,\mathcal E$ with
$\widehat{\vct e}\cdot\widehat{\vct e}^{\,*}=1$.
Here $\mathcal E$ is the scalar electric-field phasor in volts per metre.
Under the travelling-wave relation just derived,
\[
I=\frac{n\epsilon_0c}{2}|\mathcal E|^2.
\]
Define a rescaled scalar amplitude
\begin{equation}
U=\sqrt{\frac{n\epsilon_0c}{2}}\,\mathcal E
\quad\Longrightarrow\quad I=|U|^2.
\label{eq:intensity-normalized-field}
\end{equation}
The constant is absorbed into the \emph{amplitude through its square root}.
This $U$ has units $\mathrm{W}^{1/2}/\mathrm m$ if $I$ is in
$\mathrm{W/m^2}$; it is no longer the numerical electric field in
$\mathrm{V/m}$.
Multiplying by a spatially constant factor leaves the homogeneous
Helmholtz equation unchanged. In a scalar, paraxial imaging model with
a common polarization and nearly common propagation direction, this
normalization is why we may add the scalar fields and then take
$|U|^2$ to obtain intensity.

One may also divide intensity by a fixed reference intensity and use
dimensionless fields. Neither rescaling changes the interference or
edge-position argument, provided the threshold and dose use the same
normalization. For absorbing media, standing waves, or fields with
substantially different propagation directions, return to the full
flux expression in Eq.~\eqref{eq:poynting-average}; the simple
real-$n$ plane-wave prefactor is not a universal replacement for it.

\section{Optical path is not just distance}
Along a geometrical trajectory of arc length $\ell$,
\begin{equation}
\mathcal L=\int n(\vct r)\,\dd\ell.
\label{eq:opl}
\end{equation}
A uniform slab of thickness $t$ and index $n$ has optical path $nt$.
Phase accumulated along a ray is $k_0\mathcal L$ in the
$\ee^{-\ii\omega t}$ convention. Comparing actual and reference optical
paths gives OPD $W$ and phase error
\begin{equation}
\phi_W=k_0 W=\frac{2\pi W}{\lambda_0}.
\label{eq:phaseOPD}
\end{equation}
A reference must be specified: usually an ideal spherical wave converging
to a chosen image point. Moving that reference point changes the
reported tilt and defocus, even when the physical system is unchanged.

\worked{The wavelength comparison in the Introduction}
For fixed physical RMS OPD of 1.5 nm,
\[
\frac{1.5}{193}=0.0077720\ {\rm waves}
=7.7720\,\mathrm{m}\lambda,\qquad
\frac{1.5}{13.5}=0.111111\ {\rm waves}
=111.111\,\mathrm{m}\lambda.
\]
The ratio is $193/13.5=14.2963$. The length error has not grown.
Its size measured in wavelengths, and therefore its phase error, has grown.
This is the precise meaning of the paper's inverse-wavelength statement.

\section{How rays emerge from wave theory}
Write $U=A(\vct r)\ee^{\ii k_0S(\vct r)}$, where $S$ has length units.
Substitute into the Helmholtz equation:
\[
\nabla^2 A+2\ii k_0\nabla A\cdot\nabla S
+\ii k_0A\nabla^2S-k_0^2A|\nabla S|^2+n^2k_0^2A=0.
\]
When amplitude and index vary slowly on the wavelength scale, the leading
real terms require
\begin{equation}
|\nabla S|^2=n^2.
\label{eq:eikonal}
\end{equation}
This is the eikonal equation. The unit ray direction is
$\vct d=\nabla S/n$. Rays are normal to wavefronts.
The next terms give amplitude transport; ray geometry alone does not
determine interference.

For nearby actual and reference eikonals, $S=S_0+W$, the first-order
transverse direction-cosine difference is
\begin{equation}
\delta\vct d_\perp\simeq\frac{1}{n}\nabla_\perp W.
\label{eq:raygradient}
\end{equation}
Over a paraxial propagation distance $L$, the corresponding displacement
is approximately $L\nabla_\perp W/n$. This is a local geometrical statement.
A finite-aperture lithographic image still requires field summation and
source averaging.

\careful{A coefficient's sign depends on the OPD definition, time convention,
pupil axes, and image inversion. This book uses pupil phase
$\exp(+\ii2\pi W/\lambda)$ and Eq.~\eqref{eq:fourierpair}.
In transfer-function formulas the normalized pupil labels spatial
frequency, $\vct u=\vct f/f_c$; it need not have the orientation of the
physical exit-pupil coordinates in a ray drawing. Image-shift signs below
follow the declared transfer-function convention.}

\section{What a surface error does at a mirror}
For a mirror displaced a small normal distance $h$, the incident and
reflected segments each change optical path by approximately
$h\cos\alpha$, with $\alpha$ measured from the surface normal.
In vacuum and for fixed local ray geometry,
\begin{equation}
|\delta W|\simeq2|h|\cos\alpha.
\label{eq:mirror}
\end{equation}
The sign depends on positive surface displacement. At normal incidence,
a 0.1 nm height error gives about 0.2 nm OPD. Wavefront tolerance and
surface-height tolerance must not be interchanged.

This formula excludes coating phase, changed intersections, multiple-mirror
coordinate mappings, and polarization. Those belong in a full sensitivity
model. A lens-thickness perturbation in air often gives approximately
$(n-1)\delta t$ OPD rather than the mirror factor of two.

\section{Polarization and the limits of a scalar field}
The electric field is a vector perpendicular to the propagation direction
for a plane wave in an isotropic medium. The interference cross term is
$2\Real(\vct U_1\cdot\vct U_2^*)$. Orthogonal polarizations therefore have
no such cross term in total intensity. Reflection gives different
amplitudes and phases to fields perpendicular and parallel to the
incidence plane.

Scalar optics is useful when these vector effects can be separated or
neglected. At an EUV mask with finite-height absorbing features,
wavelength-scale layers, oblique illumination, and polarization-dependent
scattering, a full electromagnetic calculation may be necessary.
Rigorous simulation refers to a more complete numerical treatment of the
specified model; it does not eliminate missing inputs or manufacturing
and resist uncertainty.

\chapter{Pupils, diffraction, resolution, and focus}\label{ch:diffraction}
\reading{Saleh--Teich \S\S4.1--4.4, printed pp.~105--137.
Start at PDF pp.~127, 138, 143, and 149.
Born--Wolf \S8.3, printed p.~375 (PDF p.~413), and \S8.5.2,
printed p.~395; the Airy formula is on p.~396 (PDF p.~434).}

\section{Aperture stops, pupils, and numerical aperture}
An aperture stop limits the cone of rays that an optical system accepts
from an object point. The entrance pupil is the image of that stop seen
through the optics before it; the exit pupil is its image through the
optics after it. A pupil can therefore be a virtual image rather than a
physical hole at the plane where it is drawn.

For a rotationally symmetric image-side cone of half-angle $\alpha_{\max}$,
\begin{equation}
\NA=n\sin\alpha_{\max}.
\label{eq:NA}
\end{equation}
In vacuum $n=1$. Thus $\NA=0.33$ corresponds to
$\alpha_{\max}=\arcsin(0.33)\simeq19.27^\circ$.
The chief ray goes through the centre of the aperture stop; marginal rays
reach its edge. These are different rays.
The mask chief-ray angle of $6^\circ$ in Table 1 is not the same angle as
the wafer-side numerical-aperture cone.

In a paraxial system with pupil radius $a$ and effective image distance $f$,
$\NA\simeq a/f$ in air. The exact sine definition is the basic one.
The optical invariant gives approximately
$\NA_{\rm object}\simeq |M|\NA_{\rm image}$ for equal media and lateral
magnification $M$. A 4:1 reduction has $|M|=1/4$; object and image NAs
must not be silently interchanged.

\section{Propagation as a sum of angular components}
Insert a transverse Fourier component
$\exp[\ii2\pi(f_xx+f_yy)]\exp(\ii k_z z)$ into the Helmholtz equation
in vacuum. The result is
\[
k_z^2=k_0^2-(2\pi)^2(f_x^2+f_y^2).
\]
Each propagating component therefore acquires
\begin{equation}
H_z(\vct f)=
\exp\!\left[\ii k_0z\sqrt{1-\lambda^2|\vct f|^2}\right].
\label{eq:angularprop}
\end{equation}
For $|\vct f|>1/\lambda$, the longitudinal wave number is imaginary and
the component is evanescent. The far-field projector does not collect
arbitrarily fine near-field detail.

When $\lambda|\vct f|\ll1$,
\begin{equation}
H_z(\vct f)\simeq
\ee^{\ii k_0z}\ee^{-\ii\pi\lambda z|\vct f|^2}.
\label{eq:fresneltransfer}
\end{equation}
Inverse transformation yields the Fresnel propagation integral
\begin{equation}
U(x,y;z)=\frac{\ee^{\ii k_0z}}{\ii\lambda z}
\iint U(x',y';0)
\exp\!\left[\frac{\ii\pi}{\lambda z}
\{(x-x')^2+(y-y')^2\}\right]\dd x'\dd y'.
\label{eq:fresnel}
\end{equation}
One way to obtain the same phase is to expand the distance
$R=\sqrt{z^2+(x-x')^2+(y-y')^2}$ as
$z+\{(x-x')^2+(y-y')^2\}/(2z)$.
The amplitude $1/R$ is approximated by $1/z$.
These are two approximations, and the phase approximation is often
the stricter one because it is multiplied by the large wave number.

If the quadratic source-coordinate phase is negligible,
$a^2/(\lambda z)\ll1$ for characteristic aperture radius $a$,
Eq.~\eqref{eq:fresnel} becomes a Fourier transform:
\[
U(x,y;z)\propto
\widetilde U\!\left(\frac{x}{\lambda z},\frac{y}{\lambda z}\right)
\]
apart from a known output quadratic phase. This is Fraunhofer diffraction.
The dimensionless ratio $N_F=a^2/(\lambda z)$ is the Fresnel number.
A focusing optic can cancel the required quadratic phase, giving a
Fourier-related focal plane without requiring a physically enormous
propagation distance.

\section{The pupil function and the point-spread function}
Define a scalar coherent transfer function
\begin{equation}
P(\vct f)=A(\vct f)\,
\mathbf1_{|\vct f|\leq f_c}\,
\exp\!\left[\frac{\ii2\pi}{\lambda}W(\vct f/f_c)\right],
\qquad f_c=\frac{\NA}{\lambda}.
\label{eq:pupil}
\end{equation}
Here $\mathbf1$ equals one inside the admitted disk and zero outside,
$A$ describes amplitude transmission or reflection, and $W$ is the
phase-equivalent OPD. The normalized pupil coordinate is
$\vct u=\vct f/f_c=(\rho\cos\theta,\rho\sin\theta)$.

For coherent illumination of a thin scalar object with complex
transmittance or reflectance $t$,
\begin{equation}
U_{\rm image}=\FT^{-1}[\widetilde t\,P]=t*h,\qquad
h=\FT^{-1}[P].
\label{eq:coherentimage}
\end{equation}
$h$ is the amplitude point-spread function (PSF).
The intensity PSF is $|h|^2$. A point object contains all transverse
spatial frequencies; the pupil selects the part used to reconstruct it.
This is why a finite aperture produces a finite spot.

For a clear circular pupil with no aberration, let $r=\sqrt{x^2+y^2}$.
Angular integration gives a Bessel function:
\[
h(r)=2\pi\int_0^{f_c}fJ_0(2\pi fr)\dd f
=\frac{f_cJ_1(2\pi f_cr)}{r}.
\]
Since $h(0)=\pi f_c^2$, the normalized intensity is
\begin{equation}
\frac{I(r)}{I(0)}
=\left[\frac{2J_1(v)}{v}\right]^2,\qquad
v=2\pi f_cr.
\label{eq:airy}
\end{equation}
The first zero of $J_1$ is $v\simeq3.8317$, so
\begin{equation}
r_{\rm first\ zero}=0.6098\,\frac{\lambda}{\NA}.
\end{equation}
This is the Airy pattern. Its rings redistribute energy rather than
indicating geometric rays bouncing around the image.

\section{Coherent and incoherent resolution are different questions}
For coherent illumination, amplitudes from object points interfere before
we square the field. For mutually incoherent object points, cross terms
average out and object \emph{intensity} is convolved with $|h|^2$.
The incoherent optical transfer function is consequently a pupil
autocorrelation:
\begin{equation}
\mathrm{OTF}(\vct q)=
\frac{\int P(\vct f+\vct q/2)P^*(\vct f-\vct q/2)\dd^2f}
{\int |P(\vct f)|^2\dd^2f}.
\label{eq:otf}
\end{equation}
Two disks of radius $f_c$ cease to overlap when their separation exceeds
$2f_c$. This explains the incoherent intensity cutoff $2\NA/\lambda$,
whereas the on-axis coherent field cutoff is $\NA/\lambda$.
Partial coherence, used in lithography, requires the source-dependent
formulation developed later.

The Rayleigh two-point separation $0.61\lambda/\NA$ is a criterion for
distinguishing point-source images. It is not a universal minimum
thresholded hole diameter. Periodic patterns, oblique illumination,
resist thresholds, and proximity correction pose a different problem.
Do not conclude that the paper's 14 nm contacts are impossible merely
because $0.61(13.5)/0.33\simeq24.95\nm$.

\section{Resolution and depth-of-focus scaling}
Lithography often summarizes a process-dependent minimum feature scale as
\begin{equation}
\mathrm{CD}\sim k_1\frac{\lambda}{\NA},
\label{eq:k1}
\end{equation}
where $k_1$ absorbs pattern, illumination, mask, resist, and process choices.
For 14 nm, 13.5 nm wavelength, and 0.33 NA, the numerical ratio is
$k_1=14(0.33)/13.5\simeq0.3422$.
This calculation characterizes the geometry; it does not prove a process
window or assign the same $k_1$ to every pitch.

Defocus adds a relative phase from Eq.~\eqref{eq:angularprop}.
Subtract the on-axis piston:
\begin{align}
\phi_z(\rho)
&=\frac{2\pi z}{\lambda}
\left[\sqrt{1-\NA^2\rho^2}-1\right]\\
&\simeq-\frac{\pi z\NA^2}{\lambda}\rho^2.
\label{eq:defocusphase}
\end{align}
When the pupil-edge phase becomes a fixed tolerable number, the
corresponding defocus scales as $\lambda/\NA^2$:
\begin{equation}
\mathrm{DOF}\sim k_2\frac{\lambda}{\NA^2}.
\end{equation}
Whether DOF denotes a half-range or a full range changes the convention
for $k_2$. The paper specifies sampled focus values, not a universal $k_2$.

At its conditions,
\[
\lambda/\NA=40.91\nm,\qquad
\lambda/\NA^2=123.97\nm.
\]
Increasing NA improves lateral resolution but generally makes focus more
sensitive. Reducing wavelength improves resolution but also reduces this
focus scale at fixed NA. Reducing NA has the opposite resolution effect.
The Introduction's combined discussion of smaller wavelength and NA
should be read with these two distinct scalings in mind.

\section{A quantitative focus check}
From Eq.~\eqref{eq:defocusphase},
\[
W_z(\rho)\simeq-\frac{z\NA^2}{2}\rho^2.
\]
Remove its mean and use $\Var(\rho^2)=1/12$:
\begin{equation}
\sigma_{W,z}=\frac{|z|\NA^2}{4\sqrt3}.
\label{eq:defocusRMS}
\end{equation}
At $z=15\nm$, $\sigma_{W,z}=0.235775\nm
=17.4648\,\mathrm{m}\lambda$.
The pupil-edge-to-centre OPD is $0.81675\nm$.
These two values are different because RMS, peak-to-valley, and
edge-to-centre error are different statistics.

The exact angular formula gives an edge OPD
$15[1-\sqrt{1-0.33^2}]\simeq0.8403\nm$.
The paraxial value is a useful approximation, not an exact
high-angle propagation result.

\chapter{Why EUV needs mirrors and a three-dimensional mask}\label{ch:euv}
\paperloc{Fig.~1, Table 1, Fig.~3, and the mask discussion in Section 2A.}
\reading{Born--Wolf \S1.5, printed p.~36 (PDF p.~65), and
\S1.6, p.~51 (PDF p.~80); its characteristic-matrix discussion begins
at p.~55 (PDF p.~84). Saleh--Teich \S6.2, p.~209 (PDF p.~231).}

\section{The optical path through an EUV exposure system}
The useful sequence is source, collector and illumination optics,
reflective patterned mask, projection mirrors, and resist-coated wafer.
Illumination optics control the angular distribution incident on the
mask. Projection optics determine which mask-diffracted components reach
the wafer and with what amplitude and phase.

At 13.5 nm, strong absorption makes ordinary transmissive lens trains
and propagation through air unsuitable. EUV scanners use vacuum and
multilayer mirrors. ASML's NXE:3400C documentation specifies reflective
4:1 reduction optics and 0.33 NA \cite{asml}.
The paper's Fig.~1 caption says NXT3400; the relevant EUV family is
NXE:3400. This naming issue does not alter the optical calculations.

A geometrical mirror's reflection direction is largely independent of
wavelength, so an all-reflective design avoids ordinary refractive
chromatic aberration. Its coatings still have wavelength-dependent
reflectance and phase. Mirrors also absorb energy, heat, deform, and move.
Thus reflective does not mean perfectly achromatic in every optical
property or immune to thermal aberration.

\section{Complex refractive index and absorption}
With time convention $\ee^{-\ii\omega t}$, write a passive EUV material's
index as $\widetilde n=1-\delta+\ii\beta$, where $\beta>0$.
A normally propagating field behaves as
\[
U(z)=U(0)\ee^{\ii k_0(1-\delta)z}\ee^{-k_0\beta z}.
\]
Hence
\begin{equation}
I(z)=I(0)\ee^{-\mu_{\rm abs}z},\qquad
\mu_{\rm abs}=\frac{4\pi\beta}{\lambda}.
\end{equation}
The real part affects phase; the imaginary part attenuates amplitude.
Material optical constants therefore influence both image contrast and
phase. Treating an absorber as merely a perfectly black two-dimensional
painted region omits this physics.

\section{Reflection at an interface}
Continuity of tangential electric and magnetic fields at an interface
gives, for a chosen polarization, equations of the form
\[
1+r=t,\qquad Y_0(1-r)=Y_1t.
\]
Here $Y_j$ is the appropriate optical admittance for that polarization.
Solving,
\begin{equation}
r_{01}=\frac{Y_0-Y_1}{Y_0+Y_1},\qquad
t_{01}=\frac{2Y_0}{Y_0+Y_1}.
\label{eq:fresnelr}
\end{equation}
For electric field perpendicular to the incidence plane,
$Y_j$ is proportional to $\widetilde n_j\cos\theta_j$.
For the tangential electric-field convention of parallel polarization
it is proportional to $\widetilde n_j/\cos\theta_j$.
Different polarization amplitude conventions can reverse a reflected
amplitude sign; physical intensities and phase comparisons must use
one convention consistently.

Snell's relation is
$\widetilde n_0\sin\theta_0=\widetilde n_1\sin\theta_1$.
For absorbing layers the transmitted angle and longitudinal wave number
are generally complex. The physical branch decays into the medium.

\section{A thin film: summing all internal reflections}
Consider a layer 1 of thickness $d$ between media 0 and 2.
The phase for one transit is
$\delta_1=k_0\widetilde n_1d\cos\theta_1$.
The first reflection is $r_{01}$.
The first wave returning from the lower interface contributes
$t_{01}t_{10}r_{12}\ee^{2\ii\delta_1}$.
Each additional round trip multiplies this by
$r_{10}r_{12}\ee^{2\ii\delta_1}$.
Summing the geometric series,
\[
r_{\rm film}
=r_{01}+
\frac{t_{01}t_{10}r_{12}\ee^{2\ii\delta_1}}
{1-r_{10}r_{12}\ee^{2\ii\delta_1}}.
\]
Using $r_{10}=-r_{01}$ and
$t_{01}t_{10}=1-r_{01}^2$ gives
\begin{equation}
r_{\rm film}=
\frac{r_{01}+r_{12}\ee^{2\ii\delta_1}}
{1+r_{01}r_{12}\ee^{2\ii\delta_1}}.
\label{eq:thinfilm}
\end{equation}
Replace the lower-interface reflection by the effective reflection of
all deeper layers and repeat recursively. This constructs a multilayer
response. The equivalent transfer-matrix method propagates tangential
fields through each layer and multiplies the matrices in physical order.
The reflected intensity is $|r_{\rm stack}|^2$ in the incident medium,
and the reflected phase is $\arg r_{\rm stack}$.

This derivation explains why a small thickness error can change
reflectance and phase simultaneously. It also explains why a multilayer
reflector cannot be represented by surface height alone.

\section{The Mo/Si period and the Bragg estimate}
Reflections from successive layers can reinforce when their round-trip
phase difference is approximately $2\pi$. Neglecting refractive
corrections for an estimate, a repeated period $d$ at angle $\alpha$
from the normal obeys
\begin{equation}
2d\cos\alpha\simeq m\lambda.
\label{eq:bragg}
\end{equation}
This is a design estimate; the full result follows from the complex
multilayer calculation above, including interfaces, absorption, and
termination.

The paper gives Mo thickness $0.725\nm$, Si thickness $1\nm$, and
absorber height $15\nm$, explicitly in \emph{wafer scale}. Under a 4:1
scale conversion these correspond to physical mask values
$2.9\nm$, $4\nm$, and $60\nm$.
The inferred period is $6.9\nm$, and
\[
2(6.9)\cos6^\circ=13.7244\nm,
\]
which has the expected scale for 13.5 nm multilayer reflection.
This is a consistency check based on a 4:1 interpretation, not a
recovered complete mask prescription.

\careful{Do not put the wafer-scaled 1.725 nm period directly into a
physical 13.5 nm mask-stack solver. Lateral pattern dimensions can be
reported at wafer scale for convenience, but a Maxwell calculation must
use a consistent physical mask geometry, wavelength, and reduction mapping.
The paper does not provide all stack details needed for exact reconstruction.}

\section{Finite absorber height breaks the ideal thin-mask picture}
An obliquely incident ray traversing absorber height $h$ has a lateral
geometrical displacement $h\tan\alpha$. With the inferred
$h=60\nm$ and $\alpha=6^\circ$, this is about $6.306\nm$ on the mask,
or $1.577\nm$ after 4:1 reduction.
This shows why an apparently small angle matters on nanometre scales.
It is not a prediction of the paper's PS: coherent diffraction, material
phase, double-pass geometry, and contour formation also contribute.

Different diffraction orders interact with the finite absorber and
multilayer differently. Their complex amplitudes can depend on source
angle, polarization, hole orientation, and pitch. This produces
asymmetric images, orientation-dependent CDs, and focus-dependent
placement. A square mask opening can therefore have unequal printed
horizontal and vertical dimensions.

The mask chief-ray angle describes the central illumination geometry.
The quasar distribution describes angular offsets around that geometry.
Off-axis source shaping and a nonzero chief-ray angle are distinct ideas;
one should not infer that the chief ray itself supplies every resolution
benefit attributed to source shaping.

\section{What a rigorous mask calculation has to do}
A representative electromagnetic workflow is:
\begin{enumerate}
\item Specify each material's complex optical constants, layer thickness,
feature shape, sidewall angle, period, source angle, and polarization.
\item Solve the frequency-domain Maxwell boundary-value problem. Methods
include Fourier modal techniques and finite-element discretizations.
\item Obtain the complex vector amplitudes of reflected diffraction orders.
\item Propagate those orders through the projection system, including pupil
phase, amplitude, and defocus.
\item Sum coherent orders for one source point, average intensities over
mutually incoherent source points, and extract image or resist contours.
\end{enumerate}
For a periodic permittivity
$\epsilon(x,z)=\sum_g\epsilon_g(z)\ee^{\ii gx}$, expanding the field in
the same harmonics turns multiplication by $\epsilon$ into convolution
of Fourier coefficients. The diffraction orders become coupled equations,
not independent rays. Their boundary conditions enforce the material
interfaces. This is the basic reason rigorous diffraction is more
computationally expensive than a binary thin-mask Fourier transform.

The paper identifies a commercial Synopsys simulator but does not give
enough settings to determine its exact electromagnetic solver, numerical
convergence, material database, or resist extraction procedure.
Those items remain unknown rather than being supplied by the word rigorous.

\part{From wavefront errors to printed patterns}
\chapter{Wavefront aberrations and what RMS does not tell you}\label{ch:aberration}
\paperloc{Introduction, Section 2, and the motivation for an imaging-based budget.}
\reading{Born--Wolf \S5.1, printed pp.~203--206 (begin at PDF p.~232);
\S9.1.2, p.~462 (PDF p.~557); \S9.1.3, pp.~463--464
(PDF pp.~556--555); \S9.3, pp.~468--469 (PDF pp.~551--550).
The backward PDF order is real.}

\section{The reference sphere and the aberration function}
An ideal converging wave reaches its intended image point with the same
phase contribution from every admitted pupil location. On a suitable
reference sphere centered at that image point, its phase is constant.
The actual wave has residual phase variation $2\pi W/\lambda$.
The real-valued function $W(\rho,\theta)$ is the wave aberration.

The coordinate $(\rho,\theta)$ identifies a point in the pupil, not
a point on the wafer. A pupil map is therefore not a picture of an
aberrated contact hole. To get that picture, multiply the pupil field by
the phase factor and perform the diffraction calculation.

At a different field point on the wafer, the same optical system may
have a different pupil wavefront:
\[
W=W(\rho,\theta;\vct H),
\]
where $\vct H$ denotes image-field position. A set of coefficients for
one field point is not automatically valid across an entire exposure
field. Likewise, scan time and thermal state may add dependencies.

\section{The familiar named aberrations}
For a centred rotationally symmetric system, a low-order expansion
at field height $H$ contains terms schematically of the form
\begin{align}
W={}&A\rho^4
+B H\rho^3\cos\theta
+C H^2\rho^2\cos^2\theta\\
&+D H^2\rho^2
+E H^3\rho\cos\theta+\cdots.
\label{eq:seidel}
\end{align}
The coefficients and axes depend on convention, but the shapes reveal
the standard families:
\begin{itemize}
\item \textbf{Spherical aberration:} different pupil radii have different
focal behavior; the wavefront has a quartic radial contribution.
\item \textbf{Coma:} an off-axis asymmetric wavefront component, often
producing a lopsided point image.
\item \textbf{Astigmatism:} different meridional directions focus differently.
At a suitable intermediate focus, its wavefront has a
$\rho^2\cos2\theta$ or $\rho^2\sin2\theta$ component.
\item \textbf{Field curvature:} best focus varies across the image field.
\item \textbf{Distortion:} field-dependent placement error; a point may
remain locally sharp while its position is incorrect.
\end{itemize}
Piston, tilt, and defocus specify a constant phase, a linear pupil phase,
and a quadratic radial phase. Whether they are called aberrations or
alignment freedoms depends on the measurement purpose. The paper includes
low-order coefficients because image placement and focus behavior matter.

Equation \eqref{eq:seidel} is an illustrative classical expansion, not the
prescription of the paper's off-axis multi-mirror projector. A real EUV
system need not have the centred-system symmetries used to derive this form.

\section{Piston does not change intensity}
If $W$ is replaced by $W+c_0$, then every field contribution receives the
same factor:
\[
U_{\rm image}\mapsto\ee^{\ii2\pi c_0/\lambda}U_{\rm image}.
\]
The intensity is unchanged because
$|\ee^{\ii2\pi c_0/\lambda}|^2=1$.
Consequently intensity-derived CD, PS, and PVB cannot recover global
piston. This is an exact invariance within the model, not a failure of
neural-network training.

\section{Tilt moves the image}
Consider a linear OPD in the normalized frequency pupil,
$W=t_xu+t_yv$. Then
\[
P(\vct f)=P_0(\vct f)
\exp\!\left[\ii2\pi\left(\frac{t_x}{\NA}f_x
+\frac{t_y}{\NA}f_y\right)\right].
\]
Substituting in the inverse Fourier integral shows
\begin{equation}
U(\vct x)=U_0\!\left(\vct x+\frac{\vct t}{\NA}\right),
\qquad
\Delta\vct x=-\frac{\vct t}{\NA}.
\label{eq:tilt}
\end{equation}
Thus ideal pure tilt shifts a complete image without changing its shape.
An image registered back to its own centre has the same CD and contrast.
An unregistered image has a nonzero placement error.

The sign follows our transfer-function convention. The measurable
magnitude and the distinction between shape and placement are the
important points. When metrology removes tilt before quoting RMS,
that RMS cannot subsequently bound the removed placement error.

\section{Defocus changes the reference image plane}
Equation \eqref{eq:defocusphase} gives a quadratic pupil phase when the
observation plane moves. Equivalently, a quadratic aberration can be partly
compensated by changing focus. A report of best-focus RMS may therefore
exclude a defocus term that still matters at a fixed wafer plane.

The paper varies the focus setting by $\pm15\nm$ while also changing
aberration coefficients. To reconstruct this experiment one must know
whether the Zernike coefficients are defined at a fixed reference focus,
whether refocusing is permitted, and whether focus is reapplied separately
for every candidate. The publication does not fully specify this convention.

\section{Deriving the small-aberration Strehl relation}
For a clear, uniformly illuminated pupil and an on-axis point source,
the normalized field at the chosen reference image point is
\[
\frac{U(0)}{U_{\rm ideal}(0)}
=\avg{\ee^{\ii\phi}},\qquad \phi=\frac{2\pi W}{\lambda}.
\]
Expand the exponential:
\[
\avg{\ee^{\ii\phi}}\simeq
1+\ii\avg{\phi}-\frac12\avg{\phi^2}.
\]
Multiplying by its conjugate and retaining terms through second order,
\begin{equation}
S_{\rm ref}\simeq1-\avg{\phi^2}+\avg{\phi}^2
=1-\left(\frac{2\pi\sigma_W}{\lambda}\right)^2.
\label{eq:strehl}
\end{equation}
The exponential form often used for small wavefront errors is
\begin{equation}
S\simeq\exp[-(2\pi\sigma_W/\lambda)^2].
\label{eq:marechal}
\end{equation}
It shares the leading expansion. It is also the exact squared modulus
of the characteristic function for a zero-mean Gaussian phase distribution
under the corresponding average. An arbitrary deterministic aberration
does not make that exponential exact.

The usual Strehl ratio compares peak intensities under consistent total
power and aperture conditions. Equation \eqref{eq:strehl} initially evaluates
a chosen reference point. To discuss the optimized peak, remove or optimize
the appropriate tilt and focus and specify that choice.
With a nonuniform pupil, the amplitude weighting in the coherent field
average must also be treated explicitly.

\worked{Interpreting a $\lambda/50$ RMS requirement}
At 13.5 nm, $\lambda/50=0.27\nm=20\,\mathrm{m}\lambda$.
If this is piston-removed RMS in a suitable small-aberration setting,
\[
\sigma_\phi=2\pi/50=0.125664,\qquad
S\simeq\ee^{-(2\pi/50)^2}=0.98433.
\]
A high Strehl ratio describes concentration of a point image. It does
not guarantee a subnanometre PS for every periodic contact image.
Pure tilt illustrates the distinction particularly clearly.

\section{Equal RMS is not equal lithographic performance}
RMS measures a mean-square phase error over the pupil. A periodic mask
uses particular sets of pupil locations and their phase differences.
Two maps with identical RMS can give different phase differences at
those locations. Their CD, PS, and process window can therefore differ.

This is the physical motivation for the paper. It does not imply that
RMS is useless. RMS remains a compact measure of optical quality and a
useful control variable. It means that RMS alone is not a complete
description of a pattern-specific imaging requirement.

\chapter{Zernike polynomials from the beginning}\label{ch:zernike}
\paperloc{Equation (1), Table 1, and the coefficient axes in Figs.~5, 9, 11, and 12.}
\reading{Born--Wolf \S9.2.1, printed pp.~464--466
(PDF pp.~555, 554, 553), including its radial-polynomial table.
Its normalization differs from the RMS-normalized real basis used here.}

\section{Why use orthogonal functions?}
Suppose we tried to describe $W$ with $1,\rho^2,\rho^4,\ldots$.
These functions overlap under disk averaging:
$\avg{\rho^2}=1/2$ and $\avg{\rho^2\rho^4}=1/4$, not zero.
Changing one fitted coefficient can compensate another.
Orthogonal functions separate those overlapping contributions.

Define the inner product
\[
(f,g)=\avg{fg}
=\frac1\pi\int_0^{2\pi}\int_0^1fg\,\rho\dd\rho\dd\theta.
\]
Orthogonal means $(f,g)=0$ for different modes.
Orthonormal means each mode also has $(f,f)=1$.
The analogy is perpendicular unit axes in ordinary Euclidean geometry,
but the vectors here are functions over a disk.

\section{Constructing the first modes explicitly}
The constant mode is $Z_0^0=1$. Its RMS is one.
For a tilt, $\rho\cos\theta=u$ and $\avg{u^2}=1/4$,
so the unit-RMS tilt is
\[
Z_1^1=2\rho\cos\theta,\qquad
Z_1^{-1}=2\rho\sin\theta.
\]
To create defocus orthogonal to piston, subtract the mean from $\rho^2$:
$\rho^2-1/2$. Its mean square is
$1/3-1/4=1/12$, giving
\begin{equation}
Z_2^0=\sqrt3(2\rho^2-1).
\end{equation}
For astigmatism, angular averaging already removes piston:
\[
\avg{(\rho^2\cos2\theta)^2}
=\frac12\avg{\rho^4}=\frac16.
\]
Hence
$Z_2^2=\sqrt6\rho^2\cos2\theta$ and
$Z_2^{-2}=\sqrt6\rho^2\sin2\theta$.

For a coma-shaped mode, start from
$(\rho^3-a\rho)\cos\theta$. Remove its overlap with tilt:
\[
0=\avg{(\rho^3-a\rho)\rho\cos^2\theta}
=\frac12\left(\frac13-\frac a2\right),
\]
which gives $a=2/3$. Rescale the radial polynomial to equal one at
the rim, giving $3\rho^3-2\rho$. Its normalized real mode is
\begin{equation}
Z_3^1=\sqrt8(3\rho^3-2\rho)\cos\theta.
\label{eq:coma}
\end{equation}
The sine partner is $Z_3^{-1}$. Notice the subtracted tilt. A balanced
Zernike coma mode is not just the monomial $\rho^3\cos\theta$.

For spherical aberration, start with $\rho^4+a\rho^2+b$ and require
orthogonality to $1$ and $\rho^2$:
\[
\frac13+\frac a2+b=0,\qquad
\frac14+\frac a3+\frac b2=0.
\]
Solving gives $a=-1$, $b=1/6$. Rim normalization and unit RMS produce
\begin{equation}
Z_4^0=\sqrt5(6\rho^4-6\rho^2+1).
\label{eq:spherical}
\end{equation}
This explains why the spherical Zernike mode contains defocus and piston
terms: they make it independent of the lower radial modes.

\section{The general definition}
Let $n\geq0$ be radial degree and $m$ be an integer with
$|m|\leq n$ and $n-|m|$ even. Define
\begin{equation}
R_n^{|m|}(\rho)=
\sum_{s=0}^{(n-|m|)/2}
(-1)^s
\frac{(n-s)!}
{s!\,[(n+|m|)/2-s]!\,[(n-|m|)/2-s]!}
\rho^{n-2s}.
\label{eq:radial}
\end{equation}
The allowed parity ensures all factorials have nonnegative integer
arguments and all powers have the correct angular symmetry.
The rim value is $R_n^{|m|}(1)=1$.

Our real RMS-normalized convention is
\begin{equation}
Z_n^m=
\begin{cases}
\sqrt{n+1}\,R_n^0,&m=0,\\
\sqrt{2(n+1)}\,R_n^m\cos(m\theta),&m>0,\\
\sqrt{2(n+1)}\,R_n^{|m|}\sin(|m|\theta),&m<0.
\end{cases}
\label{eq:Zdef}
\end{equation}
The radial orthogonality identity is
\[
\int_0^1R_n^mR_{n'}^m\,\rho\dd\rho
=\frac{\delta_{nn'}}{2(n+1)}.
\]
Angular integration gives $2\pi$ for constant modes and $\pi$ for
the square of each nonconstant sine or cosine. Combining those factors
with Eq.~\eqref{eq:Zdef} gives
\begin{equation}
\avg{Z_n^mZ_{n'}^{m'}}=\delta_{nn'}\delta_{mm'}.
\label{eq:Zorth}
\end{equation}
The Kronecker delta $\delta_{ab}$ equals one if $a=b$ and zero otherwise.
This is the formal statement that the modes form independent unit axes
under this pupil-area metric.

\section{A radial table sufficient to understand low-order budgets}
Every $m>0$ entry below has both a cosine and sine partner. Apply the
normalization in Eq.~\eqref{eq:Zdef}; the table lists only radial factors.
\begin{longtable}{@{}ccp{.64\textwidth}@{}}
\toprule $n$ & $|m|$ & $R_n^{|m|}(\rho)$\\ \midrule\endhead
0&0&$1$\\
1&1&$\rho$\\
2&0&$2\rho^2-1$\\
2&2&$\rho^2$\\
3&1&$3\rho^3-2\rho$\\
3&3&$\rho^3$\\
4&0&$6\rho^4-6\rho^2+1$\\
4&2&$4\rho^4-3\rho^2$\\
4&4&$\rho^4$\\
5&1&$10\rho^5-12\rho^3+3\rho$\\
5&3&$5\rho^5-4\rho^3$\\
5&5&$\rho^5$\\
6&0&$20\rho^6-30\rho^4+12\rho^2-1$\\
6&2&$15\rho^6-20\rho^4+6\rho^2$\\
6&4&$6\rho^6-5\rho^4$\\
6&6&$\rho^6$\\
\bottomrule
\end{longtable}
There are $n+1$ real modes at degree $n$. Through degree six there are
$1+2+\cdots+7=28$ modes, including piston.
This is a complete sequence by radial degree.
It is deliberately not presented as the paper's unspecified single-index
list of 25 terms.

\figteach{zernike_atlas}{.93}{Original atlas of selected real,
RMS-normalized modes. Red and blue indicate opposite OPD signs; the same
colour scale is used throughout. The labels are $(n,m)$, not the paper's
single-index coefficient numbers.}

\section{Recovering coefficients and calculating RMS}
Expand a square-integrable wavefront as
\begin{equation}
W(\rho,\theta)=\sum_{n,m}c_n^mZ_n^m(\rho,\theta).
\end{equation}
Multiply by $Z_{n'}^{m'}$ and average. Orthogonality removes all other
terms:
\begin{equation}
c_{n'}^{m'}=\avg{WZ_{n'}^{m'}}.
\label{eq:projection}
\end{equation}
Each coefficient has the units of $W$. If $W$ is in nanometres, so is $c$.
If $W/\lambda$ is expanded, the coefficients are in waves.

Because $\avg{Z_0^0}=1$ and all other mode means are zero,
$\avg W=c_0^0$. Expanding the square and using orthogonality gives
\begin{equation}
\sigma_W^2=\sum_{(n,m)\ne(0,0)}(c_n^m)^2.
\label{eq:rss}
\end{equation}
This identity depends on the unit-RMS normalization, complete disk,
uniform area weighting, and explicit removal of piston.
Other removed modes must also be excluded. With an unnormalized
basis, the corresponding modal norm factors remain in the sum.

\worked{A coefficient is not a peak-to-valley value}
For defocus $W=cZ_2^0$, the centre value is $-\sqrt3c$ and the rim value
is $+\sqrt3c$. Thus
$\mathrm{PV}=2\sqrt3|c|$, while $\sigma_W=|c|$.
For $c=0.27\nm$, RMS is 0.27 nm but PV is about 0.9353 nm.
For the normalized tilt $W=2cu$, PV is $4|c|$.
Equal coefficients mean equal RMS for these modes, not equal PV.

\worked{Why per-mode limits are not a total RMS limit}
If 24 independent orthonormal, nonpiston modes each have magnitude
$20\,\mathrm{m}\lambda$, then
\[
\sigma_W=20\sqrt{24}\,\mathrm{m}\lambda
=97.9796\,\mathrm{m}\lambda=1.32272\nm.
\]
The input box $|c_j|\leq20\,\mathrm{m}\lambda$ therefore does not mean
total RMS is at most $\lambda/50$. If all coefficients are independent
uniform variables on that interval, then
$\E[c_j^2]=(20\,\mathrm{m}\lambda)^2/3$, and
\[
\sqrt{\E[\sigma_W^2]}=20\sqrt{24/3}
=56.5685\,\mathrm{m}\lambda.
\]
This is the square root of the expected squared RMS, not the expectation
of RMS itself. Both calculations are conditional teaching results:
the paper does not state the normalization needed to apply them directly
to its coefficient arrays.

\section{Single-index numbering is a convention, not a law}
Many programs replace $(n,m)$ by a single integer $j$.
Noll, Fringe, OSA/ANSI, and software-specific sequences differ.
Some start counting at zero, some at one; some normalize each mode to
unit RMS and some do not. Axes, sine/cosine assignment, and sign can differ.

The paper writes
\[
W(\rho,\theta)=\sum_i Z_iF_i(\rho,\theta).
\]
Here its $Z_i$ are \emph{coefficients}, while its $F_i$ are basis
functions. In this book $Z_n^m$ names a basis function and $c_n^m$ its
coefficient. These are two notation systems for the same expansion idea.
Confusing them changes the meaning of an equation.

The paper does not supply the explicit $F_i$ table or normalization.
We therefore cannot safely identify its $Z_{20}$, $Z_{21}$, or $Z_{22}$
with particular named modes. Its report that these coefficients have
narrow ranges is meaningful in its own simulator coordinates, but a
physical mode-name interpretation requires that missing dictionary.

\section{Mode orientation and rotating axes}
For a sine/cosine pair of the same $n$ and $m>0$,
\[
c_c\cos m\theta+c_s\sin m\theta
=C\cos[m(\theta-\theta_0)],
\]
where
$C=\sqrt{c_c^2+c_s^2}$ and
$\theta_0=\operatorname{atan2}(c_s,c_c)/m$.
The coefficient pair rotates when axes rotate, but $C$ stays unchanged.
Astigmatism has $m=2$, so a $45^\circ$ rotation in physical axes
interchanges the character of its cosine and sine partners.
This matters when comparing horizontal and vertical pattern sensitivities.

\section{Balanced spherical aberration as a minimization}
Suppose a surface produces raw quartic OPD $B\rho^4$. Permit a compensating
defocus $d\rho^2$ and remove piston afterward. Its variance is
\[
V(d)=B^2\Var(\rho^4)+d^2\Var(\rho^2)
+2Bd\,\Cov(\rho^4,\rho^2).
\]
Using disk averages,
\[
\Var(\rho^4)=\frac4{45},\quad
\Var(\rho^2)=\frac1{12},\quad
\Cov(\rho^4,\rho^2)=\frac14-\frac16=\frac1{12}.
\]
Differentiate:
$V'(d)=d/6+B/6=0$, giving $d=-B$.
After piston removal the balanced residual is
\[
W_{\rm bal}=B(\rho^4-\rho^2+1/6)
=\frac{B}{6\sqrt5}Z_4^0,
\]
so
\begin{equation}
\sigma_{\rm bal}=\frac{|B|}{\sqrt{180}}.
\end{equation}
The original piston-removed quartic had
$\sigma=2|B|/\sqrt{45}$, exactly four times larger.
Allowing refocus changes the budget. That change is meaningful only
if refocus is an available adjustment under the actual imaging requirement.

\section{Sampling, obscurations, and amplitude weighting}
Real measured pupils contain discrete points and may have missing regions.
For samples $\vct w$, let $B_{\ell j}=Z_j(\rho_\ell,\theta_\ell)$ and
let $Q$ be a diagonal matrix of area or uncertainty weights.
Weighted least squares minimizes
$(B\vct c-\vct w)^TQ(B\vct c-\vct w)$ and gives
\begin{equation}
(B^TQB)\vct c=B^TQ\vct w.
\label{eq:modalfit}
\end{equation}
For a masked or nonuniformly weighted pupil, $B^TQB$ is generally not
the identity. Modes orthogonal on a full disk need not remain orthogonal
under that new metric.

Pupil \emph{intensity} weighting in wavefront metrology and pupil
\emph{amplitude} weighting in coherent on-axis interference are also
different. State which quantity is weighted before interpreting RMS,
Strehl, or fitted coefficients.

\checkpoint{Why can you not square and sum an arbitrary exported list
of coefficients to get RMS? You must first know basis normalization,
pupil support and weight, units, and which modes were removed.}

\chapter{Partially coherent images of periodic contact holes}\label{ch:partial}
\paperloc{Illumination, test patterns, and the commercial imaging calculations.}
\reading{Saleh--Teich \S4.4, especially printed pp.~129--134
(PDF pp.~151--156), for coherent transfer functions.
Born--Wolf \S9.5 for coherent/incoherent imaging;
\S10.4 for mutual intensity (printed p.~507 is PDF p.~510);
Hopkins' propagation formula begins at printed p.~512 (PDF p.~505).}

\section{Why one coherent calculation is not the full image}
An extended illumination source contains many emitting regions.
The field arriving from one effective source point is coherent across
the mask for the idealized source-point calculation. Fields from distinct
source points may have rapidly varying relative phases. A resist exposure
averages over times long compared with those phase fluctuations.

Let source point $s$ produce $U_s$ with random phase $\phi_s(t)$.
Then
\[
\avg{\left|\sum_sU_s\ee^{\ii\phi_s(t)}\right|^2}_t
=\sum_s|U_s|^2+
\sum_{s\ne s'}U_sU_{s'}^*
\avg{\ee^{\ii(\phi_s-\phi_{s'})}}_t.
\]
For mutually incoherent source points the last average vanishes.
We must sum \emph{intensities} across source points, even though we
sum \emph{field amplitudes} across diffraction orders from the same point.
Confusing these two sums produces a different physical model.

More generally, the mutual intensity is
\begin{equation}
J(\vct r_1,\vct r_2)
=\avg{U(\vct r_1,t)U^*(\vct r_2,t)}_t,
\qquad
\mu_{12}=\frac{J_{12}}{\sqrt{I_1I_2}}.
\end{equation}
$|\mu_{12}|\leq1$ measures coherence; unity means full coherence at the
two points and zero means no average cross term.
If $U_{\rm out}(x)=\int h(x,r)U_{\rm in}(r)\dd r$, substitution gives
\[
I_{\rm out}(x)=
\iint h(x,r_1)h^*(x,r_2)J_{\rm in}(r_1,r_2)\dd r_1\dd r_2.
\]
The incoherent-source integral below is a convenient special form of this
general propagation rule.

\section{The source-point or Abbe imaging equation}
Let $\vct s$ be an illumination spatial frequency, expressed in
wafer-equivalent coordinates. The incident field on a thin mask is
$\ee^{\ii2\pi\vct s\cdot\vct x}$.
Multiplying the mask function $t$ by this phase shifts its spectrum:
\[
\FT[t(\vct x)\ee^{\ii2\pi\vct s\cdot\vct x}]
=\widetilde t(\vct f-\vct s).
\]
Propagate through pupil $P_z$, change the integration variable to the
mask frequency, and obtain
\begin{equation}
U_{\vct s}(\vct x)=
\ee^{\ii2\pi\vct s\cdot\vct x}
\int\widetilde t(\vct f)P_z(\vct f+\vct s)
\ee^{\ii2\pi\vct f\cdot\vct x}\dd^2f.
\end{equation}
The leading carrier has unit magnitude. For nonnegative source weights
$S$ normalized so that $\int S(\vct s)\dd^2s=1$,
\begin{equation}
I(\vct x;W,z)=
\int S(\vct s)
\left|\int\widetilde t(\vct f)P_z(\vct f+\vct s)
\ee^{\ii2\pi\vct f\cdot\vct x}\dd^2f\right|^2\dd^2s.
\label{eq:abbe}
\end{equation}
The source shifts the spectrum relative to the pupil. It does not
enlarge the physical projection aperture.

\section{The Hopkins transmission cross coefficient}
Expand the squared modulus in Eq.~\eqref{eq:abbe} and interchange
integrals:
\begin{align}
I(\vct x)
={}&\iint\widetilde t(\vct f_1)\widetilde t^*(\vct f_2)
T(\vct f_1,\vct f_2)\\
&\hspace{12mm}\times
\ee^{\ii2\pi(\vct f_1-\vct f_2)\cdot\vct x}
\dd^2f_1\dd^2f_2,\\
T(\vct f_1,\vct f_2)
={}&\int S(\vct s)
P_z(\vct f_1+\vct s)P_z^*(\vct f_2+\vct s)\dd^2s.
\label{eq:tcc}
\end{align}
$T$ is a transmission cross coefficient (TCC). It tells us how strongly
each pair of object frequencies interferes after source averaging.
Its phase contains the pupil-phase difference between the two sampled
locations. The two formulations are equivalent under their shared
thin-mask, scalar, shift-invariant assumptions.

A thick EUV mask can have a source-dependent scattering response.
Then its order amplitudes are $a_{mn}(\vct s,\mathrm{polarization})$,
and one cannot generally factor out a single source-independent
$\widetilde t$ as above. The source-point sum remains useful with the
rigorously computed amplitudes inserted separately.

\section{Diffraction coefficients of a rectangular opening}
Consider an ideal binary opening of width $a$ repeated every $p_x$
and height $b$ repeated every $p_y$. Within a period,
$t(x,y)=1$ inside the centred rectangle and zero elsewhere.
Its Fourier series is
\[
t(x,y)=\sum_{m,n}a_{mn}
\ee^{\ii2\pi(mx/p_x+ny/p_y)}.
\]
Direct integration gives
\begin{align}
a_{mn}
&=\frac{1}{p_xp_y}
\int_{-a/2}^{a/2}\ee^{-\ii2\pi mx/p_x}\dd x
\int_{-b/2}^{b/2}\ee^{-\ii2\pi ny/p_y}\dd y\\
&=\frac{ab}{p_xp_y}
\sinc\!\left(\frac{ma}{p_x}\right)
\sinc\!\left(\frac{nb}{p_y}\right),
\label{eq:rectcoeff}
\end{align}
where $\sinc q=\sin(\pi q)/(\pi q)$ and $\sinc0=1$.
The prefactor is the open-area fraction. A narrower hole spreads its
spectrum more widely, but the pupil still limits the orders that contribute.

For one source point, neglecting the irrelevant carrier,
\begin{equation}
U_{\vct s}(x,y)=\sum_{m,n}a_{mn}
P_z\!\left(\vct s+\left(\frac m{p_x},\frac n{p_y}\right)\right)
\ee^{\ii2\pi(mx/p_x+ny/p_y)}.
\label{eq:orders}
\end{equation}
The integer pair $(m,n)$ names a diffraction order.
Each admitted order is a plane-wave component of the image field.
Their relative phases and amplitudes determine the shape.

\section{Understanding the quasar illumination}
Table 1 specifies inner and outer normalized source radii 0.6 and 0.8,
angle $30^\circ$, and rotation $45^\circ$.
Figure 2 shows four off-axis source regions around diagonal directions.
Their radii are normalized to the relevant projection-pupil scale;
0.8 is not an angle of 0.8 radians.
The exact software interpretation of the angle setting is not documented.
The drawing in this book uses four annular sectors with full angular
width $30^\circ$ as an explicit teaching convention.

At pitch $38\nm$,
\[
\Delta u=\frac{1/p}{f_c}
=\frac{\lambda}{\NA p}
=\frac{13.5}{0.33(38)}\simeq1.07656.
\]
For on-axis illumination the first axial orders lie beyond a unit pupil.
Only the zero order survives in the ideal scalar square lattice,
so it cannot create the desired periodic modulation.
For a source point at radius 0.7 and azimuth $45^\circ$,
$\vct u_s=(0.49497,0.49497)$.
Now the orders $(0,0),(-1,0),(0,-1),(-1,-1)$ all lie inside the pupil.
The oblique source has made several useful orders available simultaneously.

\figteach{quasar_orders}{1.0}{Original pupil diagrams. Left: an idealized
four-sector quasar source. Right: shifted order locations for one diagonal
source point and a 38 nm square pitch. Accepted orders can interfere to
create a two-dimensional contact image. This is a thin-mask explanation,
not a reproduction of the commercial simulator.}

Increasing pitch reduces order spacing in the pupil. Different pitches
therefore sample different phases even when CD stays fixed.
Changing a rectangle's aspect ratio changes the weights $a_{mn}$ in
the two directions. This is why the authors keep many pitches and both
orientations instead of one averaged contact pattern.

\section{A two-order derivation of pattern shift}
Consider
\[
U(x)=a+b\ee^{\ii(qx+\Delta\phi)},\qquad q=2\pi/p,
\]
with real positive $a,b$. Its intensity is
\begin{equation}
I(x)=a^2+b^2+2ab\cos(qx+\Delta\phi).
\label{eq:twoorder}
\end{equation}
The maximum previously at $x=0$ is now at
\begin{equation}
\Delta x=-\frac{\Delta\phi}{q}
=-\frac{p}{2\pi}\Delta\phi.
\label{eq:fringeshift}
\end{equation}
If the differential phase is due to pupil OPD,
$\Delta\phi=2\pi(W_b-W_a)/\lambda$, then
\[
\Delta x=-\frac p\lambda(W_b-W_a).
\]
At $p=38\nm$, a differential OPD of $0.1\nm$ gives a displacement
magnitude of $0.28148\nm$ in this two-order example.
The controlling quantity is a \emph{difference at sampled pupil points},
not the pupil's global RMS.

\section{An even phase perturbation changes shape}
With symmetric side orders,
\[
U(x)=a+b\ee^{\ii\phi}\ee^{\ii qx}
+b\ee^{\ii\phi}\ee^{-\ii qx}
=a+2b\ee^{\ii\phi}\cos qx,
\]
so
\begin{equation}
I(x)=a^2+4b^2\cos^2qx+4ab\cos\phi\cos qx.
\label{eq:threeorder}
\end{equation}
The image remains even in $x$, but its intensity profile changes.
A fixed threshold therefore gives a different printed width.
Because $\cos\phi=1-\phi^2/2+\cdots$, the leading response at this
symmetric operating point is second order.
There need not be a nonzero linear sensitivity to every aberration.

\figteach{phase_to_image}{1.0}{Exact teaching examples. A differential
phase between two orders shifts fringes. A common phase on symmetric
side orders changes an even intensity profile. A real contact combines
many such interference terms and an incoherent source average.}

\section{Defocus, symmetry, and source design}
An order at normalized radius $\rho$ acquires the defocus phase
$-\pi z\NA^2\rho^2/\lambda$ in the paraxial model.
Two orders at equal radius receive the same defocus phase and preserve
their relative phase. This is one reason suitable off-axis illumination
can improve focus robustness for selected patterns.
For unequal radii, the difference is
\[
\Delta\phi_z=-\frac{\pi z\NA^2}{\lambda}
(\rho_2^2-\rho_1^2).
\]
One source design cannot generally equalize every important order pair
for all pitches and both rectangular dimensions.
Source optimization is a pattern-dependent interference problem.

\chapter{From light intensity to CD, PS, PVB, and NILS}\label{ch:metrics}
\paperloc{The six imaging indicators used throughout the paper.}

\section{An aerial image is not yet a resist contour}
The aerial image is the optical intensity in the image plane.
Exposure dose is energy per area, obtained by integrating intensity
over time. A full resist model may include photon absorption,
secondary electrons, chemical reactions, diffusion, and development.
The paper does not provide enough resist-model detail to reconstruct
that stage exactly.

To derive the basic geometry, use a declared constant-threshold model:
\begin{equation}
D\,I(x,y;z,W)=T.
\label{eq:threshold}
\end{equation}
$D$ is a relative dose multiplier and $T$ is a fixed exposure threshold
in consistent units. This is a teaching approximation. It captures how
intensity perturbations move edges without claiming a full EUV resist model.
For a bright contact opening, the printed interior is taken to be the
region above threshold; resist tone can reverse that interpretation.

\section{Critical dimension and centre position}
On a specified horizontal cut through a feature, let threshold crossings
be $x_L$ and $x_R$. Define
\begin{equation}
\mathrm{CD}_x=x_R-x_L,\qquad
x_c=\frac{x_R+x_L}{2},\qquad
\mathrm{PS}_x=x_c-x_{\rm target}.
\label{eq:CDPS}
\end{equation}
Use analogous definitions in $y$. For a nonrectangular contour, the
measurement cut, bounding box, fitted shape, or centroid definition must
be specified. A centroid computed from the whole intensity distribution
need not equal the midpoint between two threshold edges.

Small edge changes satisfy
\begin{equation}
\delta\mathrm{CD}_x=\delta x_R-\delta x_L,\qquad
\delta\mathrm{PS}_x=\frac{\delta x_R+\delta x_L}{2}.
\label{eq:edgecombine}
\end{equation}
Equal edge motion shifts the pattern. Opposite edge motion changes its
width. These two combinations explain why CD and PS provide different
information about a wavefront.

\section{Deriving the edge-displacement formula}
Let an original edge satisfy $I_0(x_e)=T$ at nominal dose one.
Perturb the image by $\delta I$ and the edge by $\delta x$:
\[
I_0(x_e+\delta x)+\delta I(x_e+\delta x)=T.
\]
Keeping first-order terms,
$I_0(x_e)+I_0'(x_e)\delta x+\delta I(x_e)=T$.
Cancel the original threshold equation:
\begin{equation}
\delta x=-\frac{\delta I(x_e)}{I_0'(x_e)}.
\label{eq:edge}
\end{equation}
The same intensity error causes a larger edge displacement when the
edge slope is shallow.
For a two-dimensional contour, normal displacement is
\[
\delta n=-\frac{\delta I}{|\nabla I_0|},
\]
with the signed normal chosen along increasing intensity.
At a vanishing slope, the linear formula fails: an opening may appear,
disappear, split, or merge. These topology changes also make learned
metric maps difficult to approximate smoothly.

\section{NILS makes the slope dimensionless}
For characteristic width $w$, define the magnitude of normalized image
log slope at an edge:
\begin{equation}
\mathrm{NILS}=w\left|\frac{\partial\ln I}{\partial x}\right|_{x_e}
=w\frac{|I'(x_e)|}{I(x_e)}.
\label{eq:nils}
\end{equation}
Different reporting conventions retain or remove the slope sign.
The dimensionless magnitude is what matters for this sensitivity argument.

A relative dose change $\epsilon=\delta D/D$ multiplies the image
by approximately $1+\epsilon$, so $\delta I=\epsilon I$.
Equation \eqref{eq:edge} gives
\begin{equation}
|\delta x|\simeq\frac{w}{\mathrm{NILS}}|\epsilon|.
\label{eq:nilsdose}
\end{equation}
For a symmetric bright feature with equal slope magnitudes at the two
edges, both edges move outward when dose increases:
\[
|\delta\mathrm{CD}|\simeq
\frac{2w}{\mathrm{NILS}}|\epsilon|.
\]
The factor of two comes from two moving edges. It must not be silently
carried into a one-edge placement formula.

\worked{A 5 percent dose perturbation}
If $w=14\nm$ and NILS is 2, one edge moves approximately
$14(0.05)/2=0.35\nm$ under a 5 percent dose change.
The full CD changes by approximately $0.70\nm$ in the symmetric example.
These assumed slopes are not measured values from the paper.
The example explains why contrast loss can enlarge process variation
even if nominal CD is adjusted back to target by changing dose.

\section{Focus-exposure matrices and process windows}
For each focus $z$ and relative dose $D$, simulate a contour and compute
$\mathrm{CD}(z,D)$, $\mathrm{PS}(z,D)$, and other required measures.
Plots of CD against focus at several doses are often called Bossung curves.
A process window is the set of focus-dose pairs satisfying \emph{all}
chosen acceptance conditions:
\[
\mathcal P=\{(z,D):\mathrm{CD}_{\min}\leq\mathrm{CD}(z,D)
\leq\mathrm{CD}_{\max},\ |\mathrm{PS}(z,D)|\leq b,\ldots\}.
\]
An attractive nominal image at $z=0,D=1$ can coexist with a very small
process window. Aberration-induced sensitivity to focus and dose is
therefore physically significant.

The paper gives $z=-15,0,+15\nm$ and dose variation $\pm5\%$.
If one uses doses $0.95,1.00,1.05$, a teaching focus-dose matrix has
nine conditions. This is a declared sampling choice here, not proof that
the authors used exactly that grid or only those dose samples.

\section{What a process variation band can mean}
For contours $C(z,D)$, their envelope is the region swept out by edge
positions over the chosen process conditions.
On a horizontal cut, define edge-band widths
\begin{equation}
B_L=\max x_L-\min x_L,\qquad
B_R=\max x_R-\min x_R.
\label{eq:pvbedges}
\end{equation}
A scalar horizontal PVB might be the larger edge-band width, a
specified average, or another software-defined contour measure.
By contrast, CD range is
\begin{equation}
B_{\rm CD}=\max(x_R-x_L)-\min(x_R-x_L).
\label{eq:cdspan}
\end{equation}
These quantities are not generally equal. For pure translation,
$B_{\rm CD}=0$ while both edge bands can be large.
For symmetric expansion/contraction, each edge band can be half
the CD range.

\careful{The paper does not give a mathematical extraction rule for
$\mathrm{PVB}_x$ and $\mathrm{PVB}_y$. Use its term as a reported
process-contour variation indicator. Do not assume without verification
that it is simply maximum CD minus minimum CD. The laboratory in this
book reports the specific edge-band widths of Eq.~\eqref{eq:pvbedges}.}

\section{Placement error, overlay, and edge placement error}
Pattern shift concerns one feature relative to its intended location.
Overlay concerns alignment between layers or exposures.
An aberration-induced shift can consume part of an overlay allowance,
but stage alignment, mask registration, wafer distortion, and other
contributors also matter.

For a one-dimensional target with left and right edges, small
edge placement errors combine CD and PS:
\begin{equation}
\mathrm{EPE}_R=\mathrm{PS}+\frac{\delta\mathrm{CD}}2,\qquad
\mathrm{EPE}_L=\mathrm{PS}-\frac{\delta\mathrm{CD}}2.
\label{eq:epe}
\end{equation}
This algebra shows why controlling only the feature centre is insufficient
to control every edge. A contact can be correctly centred but too large.
Conversely, its width can be perfect while its entire position is wrong.

\section{Why horizontal and vertical errors differ}
The mask incidence plane establishes a preferred direction. A rectangular
hole has different spectra in $x$ and $y$. A quasar source has discrete
angular structure. A wavefront mode has its own orientation.
The interplay can make one direction more sensitive than the other.

The paper's relatively large $\mathrm{PS}_x$ and small
$\mathrm{PS}_y$ should not be interpreted as a universal property of all
EUV contacts. The incidence azimuth, coordinate conventions, baseline
placement, and definition of the reported shifts would need to be known.
A small denominator in a percentage error is a separate statistical issue,
addressed in Chapter \ref{ch:errors}.

\chapter{How an aberration coefficient changes a metric}\label{ch:sensitivity}
\section{Differentiating the pupil and the field}
Let $W=\sum_jc_jZ_j$ with coefficients in length units.
Then
\begin{equation}
\frac{\partial P}{\partial c_j}
=\frac{\ii2\pi}{\lambda}Z_jP.
\label{eq:pupilDerivative}
\end{equation}
For a fixed thin-mask spectrum and one source point,
\[
\frac{\partial U_{\vct s}}{\partial c_j}
=\frac{\ii2\pi}{\lambda}
\int \widetilde t(\vct f)\,
Z_j\!\left(\frac{\vct f+\vct s}{f_c}\right)
P(\vct f+\vct s)\ee^{\ii2\pi\vct f\cdot\vct x}\dd^2f.
\]
Differentiate $I=UU^*$:
\begin{equation}
\frac{\partial I}{\partial c_j}
=2\Real\left\{U^*\frac{\partial U}{\partial c_j}\right\}.
\label{eq:intDerivative}
\end{equation}
For a partially coherent source, integrate this expression over source
weights. In a vector calculation, replace the scalar product with the
appropriate polarization sum.

The derivative is physically a measure of how a phase change in one
mode perturbs the constructive and destructive interference in the image.
It is not determined merely by how large the mode looks in a pupil plot.

\section{Differentiating a printed edge}
Apply implicit differentiation to
$I(x_e(\vct c),\vct c)=T$:
\[
\frac{\partial I}{\partial x}\frac{\partial x_e}{\partial c_j}
+\frac{\partial I}{\partial c_j}=0.
\]
Therefore
\begin{equation}
\frac{\partial x_e}{\partial c_j}
=-\frac{\partial I/\partial c_j}{\partial I/\partial x}.
\label{eq:edgeSensitivity}
\end{equation}
Combine left and right derivatives using Eq.~\eqref{eq:edgecombine}
to obtain the sensitivity of CD and PS. This supplies an explicit
physics-based path from coefficient to metric.

For $M$ measured indicators, define
\begin{equation}
J_{ij}=\left.\frac{\partial F_i}{\partial c_j}\right|_{\vct c_0},
\qquad
\delta\vct y\simeq J\,\delta\vct c.
\label{eq:sensitivity}
\end{equation}
$J$ is the imaging Jacobian. Its rows correspond to particular patterns,
orientations, and metrics; its columns correspond to modes.
If coefficients and metrics are both in nm, its entries are dimensionless.
If coefficients are in milliwaves, the entries have units
nm per milliwave.

\section{A numerical derivative when an analytic derivative is unavailable}
With a black-box simulator, a centred finite difference gives
\begin{equation}
J_{ij}\simeq
\frac{F_i(\vct c_0+h\vct e_j)-F_i(\vct c_0-h\vct e_j)}{2h}.
\label{eq:finitedifference}
\end{equation}
Here $\vct e_j$ changes only coefficient $j$.
Taylor expansion shows the leading truncation error is $O(h^2)$ for
a sufficiently smooth map.
Making $h$ extremely small is not always beneficial: numerical contour
interpolation, solver convergence error, and cancellation can dominate.
Check derivative stability over a sensible range of $h$ and ensure
the same physical feature is being measured.

\section{Interactions and nonlinear response}
The second-order expansion of indicator $i$ is
\begin{equation}
\delta y_i\simeq\sum_jJ_{ij}\delta c_j+
\frac12\sum_{j,k}H_{ijk}\delta c_j\delta c_k,
\end{equation}
where $H_{ijk}$ is a second derivative. Terms with $j\ne k$ describe
interactions between modes. Single-mode sweeps do not by themselves
measure those cross terms.

For example, if $y=c_1c_2$, varying either coefficient while keeping the
other zero produces no response. Yet both nonzero together produce a
finite response. This elementary example explains why the paper
supplements individual-mode sweeps with combined aberrations.

\section{Symmetry predicts some zero sensitivities}
If a complete imaging configuration is invariant under reflection
$x\mapsto-x$, an even phase perturbation cannot by itself generate an
odd placement change under conditions preserving that symmetry.
Likewise, a quantity may depend on a coefficient through $c^2$ near zero,
giving zero first derivative but nonzero second derivative.

These conclusions require symmetry of the whole experiment: source,
mask, pupil amplitude, aberration, focus, and metric definition.
Oblique EUV mask scattering can destroy a symmetry present in a
simplified scalar calculation. Mode names alone cannot establish
which physical direction shifts.

\section{Why a physical inverse is difficult but not impossible}
The paper says required aberrations cannot be deduced using physical
laws or equations because of complexity. The careful interpretation is
that a convenient closed-form inverse is generally unavailable.
The forward map is still governed by wave physics; its derivatives can
be computed, and iterative constrained optimization can use them.
The neural network offers an approximate, data-driven route to candidate
solutions. It does not replace Maxwell's equations or prove that
physics-based inverse design is impossible.

\checkpoint{Can a coefficient with a small derivative at zero be given
an unlimited budget? No. Higher-order effects, interactions, changes in
the active contour, and other patterns can become important away from zero.}

\part{Understanding the authors' inverse model}
\chapter{Reconstructing the paper's dataset and assumptions}\label{ch:data}
\paperloc{Section 2A, Tables 1--3, and the preprocessing in Section 2C.}

\section{The optical conditions that define this particular experiment}
An aberration-to-image map is conditional on many fixed inputs.
More explicitly, one should write
\[
\vct y=F(\vct c;\lambda,\NA,S,\mathcal M,\mathcal R,\mathcal Q),
\]
where $S$ is illumination, $\mathcal M$ is mask geometry and material
response, $\mathcal R$ is resist or contour extraction, and
$\mathcal Q$ is the set of sampled process conditions.
Changing any of these changes the map to be learned.

\begin{table}[htbp]\centering\small
\caption{The paper's Table 1, expressed with explicit meanings.
Values are reported from the publication; qualifications below are this
guide's interpretation.}
\begin{tabularx}{\textwidth}{@{}p{.31\textwidth}Y@{}}
\toprule Quantity & Reported value\\ \midrule
Vacuum wavelength & 13.5 nm\\
Projection NA & 0.33\\
Chief-ray angle at object & $6^\circ$\\
Illumination & Quasar: inner radius 0.6, outer radius 0.8,
angle $30^\circ$, rotation $45^\circ$\\
Absorbing topography & $\mathrm{Ta}_6\mathrm{N}_4$,
15 nm in wafer scale\\
Reflecting multilayer & Mo/Si, respectively 0.725 nm and 1 nm
in wafer scale\\
Focus settings & $-15,0,+15\nm$\\
Exposure-dose variation & $\pm5\%$\\
Zernike indices & $Z_2$ through $Z_{25}$\\
Coefficient sweep & $-20$ to $+20\,\mathrm{m}\lambda$,
step $1\,\mathrm{m}\lambda$\\
\bottomrule
\end{tabularx}
\end{table}

The optical constant database, complete multilayer termination,
number of periods, exact source discretization, polarization,
mask sidewalls, pupil apodization, and resist settings are not specified
well enough for a numerical reconstruction. A student should understand
their role rather than infer them from typical values.

\section{The intended pattern families in Table 2}
\begin{table}[htbp]\centering\small
\caption{Pattern specification reported in the paper's Table 2.
Each pitch list contains seven entries; rectangular pitch pairs are
read position by position.}
\begin{tabularx}{\textwidth}{@{}L{.22\textwidth}L{.14\textwidth}Y Y@{}}
\toprule Family & $(\mathrm{CD}_x,\mathrm{CD}_y)$ (nm)
& $p_x$ (nm) & $p_y$ (nm)\\ \midrule
Square vias & $(14,14)$ & 38, 40, 42, 46, 50, 80, 120
& 38, 40, 42, 46, 50, 80, 120\\
Rectangular vias & $(14,28)$ & 56, 60, 64, 70, 80, 100, 150
& 70, 74, 78, 84, 94, 114, 164\\
\bottomrule
\end{tabularx}
\end{table}
Through-pitch testing means holding a nominal feature type or size while
varying its neighbourhood spacing. This probes proximity effects and
different pupil sampling. It does not mean changing the wavelength or
continuously scanning every possible pitch.

The count suggested by two families and seven variants is 14 patterns.
Each contributes six scalar indicators:
\[
(\mathrm{CD}_x,\mathrm{CD}_y,\mathrm{PS}_x,\mathrm{PS}_y,
\mathrm{PVB}_x,\mathrm{PVB}_y).
\]
Thus $14\times6=84$ explains the reported network input width.
This is a consistent reconstruction of the feature count, not a
published list of the exact vector ordering.

\section{The conflicting plot-label table}
The paper's Table 3 assigns the following parameters to the bars labelled
REC and SQ. It is important to preserve this discrepancy:
\begin{table}[htbp]\centering\small
\caption{Paper Table 3 as printed. Every row uses CD $(14,28)$ nm for
REC and $(14,14)$ nm for SQ. These pitch assignments disagree with Table 2.}
\begin{tabular}{@{}crrrr@{}}
\toprule & \multicolumn{2}{c}{REC pitch (nm)}
& \multicolumn{2}{c}{SQ pitch (nm)}\\
Pattern & $p_x$ & $p_y$ & $p_x$ & $p_y$\\ \midrule
1&38&38&56&56\\
2&40&40&60&74\\
3&42&42&64&78\\
4&46&46&70&84\\
5&50&50&80&94\\
6&80&80&100&114\\
7&120&120&150&164\\
\bottomrule
\end{tabular}
\end{table}
The pitch families appear exchanged relative to the shape assignments
in Table 2. The first SQ pair $(56,56)$ also differs from the rectangular
pair $(56,70)$ in Table 2.
The text says REC means rectangle and SQ means square, but does not
resolve which table reflects the actual simulation layout.

When reading the graphs, retain the authors' labels and pattern numbers.
When performing this book's educational calculation, use explicitly
declared geometry. Do not claim that a chosen correction reproduces
the authors' dataset.

\section{Three groups of aberration inputs}
The stated simulation sets are:
\begin{enumerate}
\item \textbf{Individual-mode sweeps:} 41 coefficient values
$-20,-19,\ldots,+20\,\mathrm{m}\lambda$ for each of 24 indices,
with all others zero. This gives $41\times24=984$ rows.
\item \textbf{Definitive screening design:} 77 combined cases using
levels $-20,0,+20\,\mathrm{m}\lambda$.
\item \textbf{Random combinations:} 972 additional cases.
\end{enumerate}
Their nominal total is $984+77+972=2033$ before any filtering.
The paper does not give a complete split into training, validation,
and testing, nor the exact random distribution, seed, or retained row count.

The individual sweeps include the all-zero vector 24 times.
They therefore contain at most $984-23=961$ distinct vectors.
Duplicates between groups may further reduce the distinct count.
This matters in a replication: randomly splitting duplicate optical
cases between training and test sets would weaken the independence
of the test. The paper does not establish that such leakage occurred;
it simply leaves the handling undocumented.

\section{Why a screening design is used}
A local quadratic response in $K$ variables has
\[
1+K+K+\frac{K(K-1)}2
\]
coefficients: intercept, linear terms, squares, and pairwise interactions.
For $K=24$ this is $325$.
An efficient screening design seeks useful information about dominant
effects from far fewer than a full factorial set of runs, usually under
sparsity or low-order assumptions.

Three levels allow curvature to be distinguished from a simple straight
line. Carefully chosen combinations can separate some effects better
than arbitrary combinations. However, 77 runs cannot independently
estimate all 325 coefficients of an unrestricted quadratic model.
The exact design matrix and its aliasing properties are not supplied.
No specific orthogonality guarantee for these 77 rows can be reconstructed.

If every coefficient took all 41 grid values, a full 24-dimensional
grid would contain
\[
41^{24}=5.091110945\ldots\times10^{38}
\]
points. This count is a hypothetical grid comparison, not a claim that
the random group follows that grid. It illustrates why a few thousand
cases can only sample a minute portion of the possible coefficient space.

\section{What is one machine-learning sample?}
The natural forward row is
\[
\vct c^{(r)}
\stackrel{F}{\longrightarrow}
\vct y^{(r)}\in\R^{84},
\]
where the vector holds all patterns and metrics for aberration case $r$.
The inverse training row reverses the pair:
$\vct y^{(r)}$ is input, $\vct c^{(r)}$ is target.
The 14 patterns are therefore features within one optical case.
They are not automatically 14 independent realizations of the same input.

Keeping them separate is physically useful. For two patterns with
responses $\delta y_1=c$ and $\delta y_2=-c$, averaging produces zero
and loses all sensitivity to $c$. A feature vector $(c,-c)$ preserves it.
This is the paper's motivation for avoiding averages over pattern
performance during learning.

\section{The 24-versus-25 output issue}
The inclusive range $Z_2,\ldots,Z_{25}$ contains 24 coefficients.
The $984=41\times24$ count and plotted axes support 24 varied terms.
However, the prose repeatedly refers to the first 25 terms and the
reported network has 25 output neurons.

One possible explanation is inclusion of a fixed piston coordinate,
but the paper does not confirm that explanation. If piston is included
as an unconstrained output, intensity indicators cannot determine it.
A reconstruction must establish the exact output dictionary rather
than fabricate a twenty-fifth observable aberration.

\section{Filtering and scaling are part of the model}
Section 2C describes filtering poorly fitted or low-valued data and
mentions a 0.001 threshold, but its wording mixes input data and
output Zernike terms. It does not establish which rows or coordinates
were removed, the units before filtering, or whether the decision used
test information.

This ambiguity affects reproducibility and performance interpretation.
For example, removing all near-zero shift examples can improve a
percentage-error score while making predictions near the desired zero
shift less well supported. Filtering must be defined from the training
workflow and recorded; it is not just a harmless plotting choice.

\chapter{The inverse problem: existence, ambiguity, and feasible targets}\label{ch:inverse}
\section{Three separate questions}
Given desired metrics $\vct y_*$, ask:
\begin{enumerate}
\item \textbf{Existence:} is there any allowed wavefront with
$F(\vct c)=\vct y_*$, or sufficiently close?
\item \textbf{Uniqueness:} if a solution exists, is it the only one?
\item \textbf{Stability:} do small changes in desired metrics lead to
small changes in the required coefficients?
\end{enumerate}
A forward simulator may be perfectly deterministic even when the
inverse fails one or more of these tests.

\section{More indicators than coefficients does not settle uniqueness}
With about 24 varied coefficients and 84 metrics, there are more
reported numbers than unknowns. This does not imply either uniqueness
or nonuniqueness by itself. Metrics can be correlated, insensitive,
or redundant; nonlinear maps can also have multiple branches.

Locally, $J$ has 84 rows and 24 columns. If it has full column rank,
small coefficient changes are locally distinguishable by the measured
metrics. If its rank is lower, some first-order changes are invisible.
Even full local rank does not rule out two distant coefficient vectors
with the same metrics.

Piston is an exact unobservable direction if included. A simple
additional teaching example is $F(c)=c^2$: both $c=1$ and $c=-1$
give the same output. In symmetric optical settings, sign ambiguities
can arise from similar even dependence, but they should be established
for the actual source and mask rather than assumed from a mode name.

\section{The singular-value decomposition explains sensitivity}
Any real Jacobian can be written
\begin{equation}
J=U\Sigma V^T.
\label{eq:svd}
\end{equation}
The columns of $V$ are orthogonal coefficient-space directions.
Applying $J$ to direction $\vct v_j$ gives
$J\vct v_j=\sigma_j\vct u_j$.
The singular value $\sigma_j$ measures how much that coefficient
combination changes the measured outputs.

A small $\sigma_j$ means weak observability. Formally inverting this
component divides by $\sigma_j$, amplifying measurement and modelling
error. A zero value is a linear null direction.
This is a more precise concept than calling one single Zernike term
unimportant: a weak direction can be a correlated combination of many terms.

Before calculating an SVD, scale coordinates appropriately. Otherwise
replacing nm by milliwaves changes numerical singular values simply by
changing units. A useful scaled Jacobian uses coefficient allowance
scales in its columns and metric tolerance scales in its rows.

\section{Linear inverse and regularization}
For desired change $\delta\vct y$, a weighted least-squares inverse solves
\begin{equation}
\min_{\delta\vct c}
\frac12\norm{Q_y^{1/2}(J\delta\vct c-\delta\vct y)}^2
+\frac{\alpha}{2}\norm{L\delta\vct c}^2.
\label{eq:regularized}
\end{equation}
$Q_y$ weights metric errors; $L$ can penalize coefficient magnitude
or difficult-to-manufacture combinations; $\alpha\geq0$ balances the terms.
Differentiation gives
\begin{equation}
(J^TQ_yJ+\alpha L^TL)\delta\vct c=J^TQ_y\delta\vct y.
\end{equation}
With $Q_y=I$ and $L=I$, an SVD component is multiplied by
$\sigma_j/(\sigma_j^2+\alpha)$ rather than $1/\sigma_j$.
Small, unreliable directions are suppressed.
The penalty is a design preference, not additional optical data.

\section{Feasible outputs occupy a structured subset}
If $\vct c$ varies in a $K$-dimensional region, its smooth image in
84-dimensional metric space has at most $K$ local degrees of freedom.
Many metric combinations are therefore unattainable.

Being inside the minimum and maximum of each training feature is a
necessary range check, but not a sufficient physical-feasibility check.
For example, data satisfying $y_2=y_1^2$ with $-1\leq y_1\leq1$ have
coordinate ranges $y_1\in[-1,1]$ and $y_2\in[0,1]$.
The point $(0,1)$ lies inside both ranges but not on the curve.

The paper's 118 target vectors come from other simulations under the same
lithographic conditions and were not used in training or validation,
according to Section 3. That is stronger than choosing 84 coordinates
independently within a box. The paper nevertheless does not release the
vectors or their generation procedure, so their precise coverage is unknown.

\section{Why averaging inverse solutions can fail}
Suppose training pairs have identical input $y=1$ and two valid targets
$c=+1$ and $c=-1$ under $F(c)=c^2$.
A deterministic model trained by squared coefficient error minimizes
\[
L(\hat c)=\tfrac12[(\hat c-1)^2+(\hat c+1)^2]=\hat c^2+1.
\]
Its minimizer is $\hat c=0$, but $F(0)=0\ne1$.
The mean of valid inverse solutions is not necessarily valid.

More generally, squared-error regression favors the conditional mean
$\E[\vct c\mid\vct y]$, which can fall between distinct feasible branches.
A good test is therefore the forward consistency
\begin{equation}
\vct r_{\rm image}=F(G(\vct y))-\vct y,
\label{eq:forwardconsistency}
\end{equation}
not only the coefficient residual $G(\vct y)-\vct c_{\rm original}$.
The paper's forward re-simulation is an important part of its logic.
It checks whether a proposed alternative actually retains useful imaging.

\section{One output per input versus a distribution of solutions}
A deterministic feedforward network returns one coefficient vector
for one specified metric vector. Giving it 118 different target vectors
can produce 118 candidates. That does not mean the same trained network
samples all allowable wavefronts for a single fixed target.

To obtain diverse alternatives for the same target one could use multiple
optimization initializations, explicit latent inputs, or a conditional
generative model, followed by forward verification. These are possible
extensions, not methods reported in this paper.

\section{The feasible set is the real object of a budget}
Let permitted coefficient conditions be $\vct c\in\mathcal C$ and
imaging limits be $\vct y\in\mathcal Y_{\rm allowed}$.
The acceptable wavefront set is
\begin{equation}
\mathcal A=\{\vct c\in\mathcal C:F(\vct c)\in\mathcal Y_{\rm allowed}\}.
\label{eq:feasibleset}
\end{equation}
The paper identifies a finite collection of candidates in or near this set,
as judged by its simulator. The shape, volume, boundary, robustness, and
manufacturability of the entire set require further analysis.
Chapter \ref{ch:budget} develops the distinction between candidate
examples and a guaranteed tolerance region.

\chapter{The neural network and its training, step by step}\label{ch:network}
\paperloc{Sections 2B--2C, Fig.~4, and the reported network architecture.}

\section{A neuron is a weighted sum followed by a function}
Given inputs $a_1,\ldots,a_d$, a neuron forms
\[
q=\sum_{j=1}^dw_ja_j+b,\qquad a_{\rm out}=g(q).
\]
The weights $w_j$ determine which inputs contribute and with what sign;
the bias $b$ shifts the response. The activation $g$ introduces nonlinearity.
A layer of neurons applies this operation in parallel:
\begin{equation}
\vct q^{(\ell)}=A^{(\ell)}\vct a^{(\ell-1)}+\vct b^{(\ell)},
\qquad
\vct a^{(\ell)}=g_\ell(\vct q^{(\ell)}).
\label{eq:network}
\end{equation}
The function acts component by component. Repeating layers produces a
composition of nonlinear functions. With only linear activations,
all layers collapse into a single affine map, no matter how many are used.

For example, $g(q)=\tanh q$ is smooth, ranges from $-1$ to $+1$,
and has derivative $1-\tanh^2q$. A rectified linear unit is
$\max(0,q)$ and is piecewise linear. The paper does not report its
activation functions, so neither choice can be attributed to the authors.
Likewise, an ordinary unbounded linear regression output is a possible
choice, but the actual output activation is undocumented.

\section{What the network is asked to learn}
Here the input is the 84-component imaging vector, not a pupil image.
The output is a vector of Zernike coefficients, not CD or PVB.
The desired relationship is
\[
G_\Theta:\vct y\longmapsto\vct c,
\]
where $\Theta$ collects all trainable weights and biases.
During training, the known simulator pair
$(\vct y^{(r)},\vct c^{(r)})$ supplies the input and target.
During use, a new desired imaging vector supplies the input.

Calling this a backpropagation network describes how derivatives of the
training loss are computed. It does not mean light is physically
propagated backward through a Maxwell solver.
Back prediction refers to the direction of the learned inverse mapping;
backpropagation refers to the chain rule used to train it.

\section{Interpreting the reported layer widths}
The paper gives
\[
84,\ 84,\ 64,\ 64,\ 50,\ 36,\ 32,\ 25,\ 25.
\]
Read literally, the network has 84 inputs and 25 outputs. Its seven
hidden layers have widths $84,64,64,50,36,32,25$.
Figure 4 is a simplified schematic and does not depict these exact widths
or the full hidden-layer count.

A fully connected layer taking $d_{\rm in}$ inputs to $d_{\rm out}$
outputs has $d_{\rm in}d_{\rm out}$ weights and $d_{\rm out}$ biases.
The literal architecture therefore has
\[
\sum_{\ell=1}^8(d_{\ell-1}+1)d_\ell
=7140+5440+4160+3250+1836+1184+825+650
=24\,485
\]
trainable parameters. This calculation assumes the conventional dense
architecture; the publication does not provide implementation code to
confirm all details.

The parameter count exceeds the number of simulated aberration cases.
That does not mathematically prove overfitting: each case has multiple
outputs, and training structure matters. It does make independent testing,
preprocessing disclosure, and regularization choices important.
The 25-output ambiguity discussed in Chapter \ref{ch:data} remains.

\section{Why feature normalization matters}
CD may be of order tens of nanometres while $\mathrm{PS}_y$ can be a small
fraction of a nanometre. An optimizer acting on raw coordinates may assign
very different numerical scales to these quantities.
Featurewise min--max normalization to $[-1,1]$ is
\begin{equation}
\widetilde x=2\frac{x-x_{\min}}{x_{\max}-x_{\min}}-1.
\label{eq:normalization}
\end{equation}
Solving for the physical value,
\begin{equation}
x=x_{\min}+\frac{\widetilde x+1}{2}(x_{\max}-x_{\min}).
\label{eq:unnormalization}
\end{equation}
The paper names MATLAB's \texttt{mapminmax}, which performs rowwise
range mapping and supports applying stored settings to new data
\cite{mapminmax}.

The minimum and maximum must be fitted on training data and saved.
Use those same settings for validation, testing, and deployment.
Recomputing them independently for a new batch changes the meaning
of every network coordinate. Including held-out data when fitting them
lets test information enter preprocessing.

If $x_{\max}=x_{\min}$, the feature has no training variation and the
formula divides by zero. It needs explicit handling.
If a new value lies outside the training range, its normalized value
lies outside $[-1,1]$; min--max mapping does not by itself clip it.
Similarly, undoing output normalization does not impose a coefficient
bound unless the output is explicitly bounded or clipped.

\section{The training loss}
Let $N$ be the number of training samples and $K$ the number of output
coordinates actually trained. One common mean-square loss is
\begin{equation}
L(\Theta)=\frac{1}{NK}\sum_{r=1}^N\sum_{j=1}^K
\left[\widehat c_j^{(r)}-c_j^{(r)}\right]^2.
\label{eq:trainingMSE}
\end{equation}
Often the coefficients here are normalized coordinates.
The paper says training uses MSE and a goal of $10^{-8}$, but does not
supply enough preprocessing detail to turn that number into a guaranteed
physical coefficient error.

If, purely as an example, one output spans
$[-20,+20]\,\mathrm{m}\lambda$ and is mapped to $[-1,1]$, then
a normalized RMSE $10^{-4}$ corresponds to
$0.002\,\mathrm{m}\lambda$ for that coordinate.
The square root of an average training MSE does not bound each output,
does not describe held-out error, and does not bound image error after
the nonlinear forward map.

\section{Backpropagation is repeated application of the chain rule}
For clarity use a one-sample loss
$L=\tfrac12\norm{\widehat{\vct c}-\vct c}^2$.
For a linear output layer, the derivative with respect to its preactivation is
\[
\vct\delta^{(L)}=\widehat{\vct c}-\vct c.
\]
For a nonlinear output, multiply componentwise by $g_L'$.
At each earlier layer,
\begin{equation}
\vct\delta^{(\ell)}
=\left[(A^{(\ell+1)})^T\vct\delta^{(\ell+1)}\right]
\odot g_\ell'(\vct q^{(\ell)}).
\label{eq:backprop}
\end{equation}
The symbol $\odot$ means componentwise multiplication.
To understand the transpose, note that one hidden neuron influences
every neuron in the next layer; its derivative sums all those downstream
paths, weighted by the connecting weights.

The parameter derivatives are
\begin{equation}
\frac{\partial L}{\partial A^{(\ell)}}=
\vct\delta^{(\ell)}(\vct a^{(\ell-1)})^T,\qquad
\frac{\partial L}{\partial\vct b^{(\ell)}}=\vct\delta^{(\ell)}.
\label{eq:networkgrad}
\end{equation}
Average or sum over samples as required by the loss normalization.
Backpropagation computes these derivatives efficiently by reusing
intermediate derivatives, rather than perturbing every weight separately.

\worked{One hidden neuron}
Take $a=\tanh(wx+b)$, $\hat c=va+d$, and
$L=\tfrac12(\hat c-c)^2$.
With $e=\hat c-c$,
\[
\frac{\partial L}{\partial v}=ea,\quad
\frac{\partial L}{\partial d}=e,\quad
\frac{\partial L}{\partial w}=ev(1-a^2)x,\quad
\frac{\partial L}{\partial b}=ev(1-a^2).
\]
Every factor has a role: residual, downstream weight, activation slope,
and upstream input. A deep network repeats the same multiplication and
summation pattern across more paths.

\section{Deriving the Levenberg--Marquardt update}
Flatten all training residuals into a vector $\vct e(\Theta)$ and let
$J_e=\partial\vct e/\partial\Theta$.
This \emph{training Jacobian} differentiates residuals with respect to
network parameters. It is distinct from the imaging Jacobian
$\partial\vct y/\partial\vct c$.

For a small parameter change $\Delta\Theta$,
$\vct e(\Theta+\Delta\Theta)\simeq\vct e+J_e\Delta\Theta$.
Minimize the locally linearized residual with a damping penalty:
\[
\frac12\norm{\vct e+J_e\Delta\Theta}^2
+\frac{\mu}{2}\norm{\Delta\Theta}^2.
\]
Differentiate with respect to $\Delta\Theta$, set the derivative to zero,
and obtain
\begin{equation}
(J_e^TJ_e+\mu I)\Delta\Theta=-J_e^T\vct e.
\label{eq:LM}
\end{equation}
Solve this linear system; do not explicitly form a matrix inverse
in a numerical implementation.

For small $\mu$, the step approaches Gauss--Newton.
For very large $\mu$,
$\Delta\Theta\simeq-\mu^{-1}J_e^T\vct e$,
which is a small gradient-descent step.
The optimizer adjusts damping according to whether a proposed step
improves the loss. This is the mathematical reason LM can combine
rapid local convergence with caution away from a good fit.

For a squared-residual loss, the exact Hessian contains
$J_e^TJ_e+\sum_i e_i\nabla^2e_i$.
Gauss--Newton drops the second term, which is often small near a
good fit. LM adds the stabilizing diagonal term.
It does not guarantee finding a global optimum of a deep network.

The paper identifies MATLAB's LM training function. Its official
documentation distinguishes the MSE performance goal from
validation-based stopping and documents the damping update
\cite{trainlm}. These software facts confirm the interpretation of the
algorithm; the derivation above follows directly from the local
least-squares problem.

\section{Training, validation, and test have different jobs}
Training data determine the weights. Validation data guide choices such
as stopping time or architecture. Test data estimate performance after
those choices are fixed. Repeatedly changing the architecture after
looking at the same test set makes that set part of model selection.

The paper says training stops below MSE $10^{-8}$ or after six rounds
without an obvious decrease. Standard MATLAB \texttt{trainlm} uses a
six-failure default for \emph{validation} performance.
The paper's wording does not settle whether its implemented condition
was precisely this default or a different training criterion.
Record that ambiguity rather than equating training stagnation with
validation failure.

\section{What the learning theory does and does not promise}
The Introduction cites approximation results for neural networks.
Such results motivate using nonlinear networks to approximate continuous
functions under suitable conditions. They do not establish that this
particular architecture, finite dataset, and optimizer will recover a
stable inverse of an ambiguous optical map.

In particular, they do not guarantee extrapolation beyond the training
region, compliance with coefficient bounds, a manufacturable optical
system, or reliable behavior when the source or mask changes.
The paper's useful evidence is its forward-simulation validation,
not the general existence of approximation theorems.

\chapter{Errors, percentages, and what validation means}\label{ch:errors}
\paperloc{Equations (2)--(3), Figs.~5--10, and the statements about small PS values.}

\section{Root mean square error}
For predictions $\hat y_i$ and reference values $y_i$, define
$e_i=\hat y_i-y_i$. Then
\begin{equation}
\mathrm{RMSE}=\sqrt{\frac1N\sum_{i=1}^Ne_i^2}.
\label{eq:RMSE}
\end{equation}
This is the paper's Eq.~(2). Squaring prevents positive and negative
errors from cancelling. Taking the square root restores the units.
RMSE of a CD in nm is in nm. RMSE of a coefficient in waves is in waves.
The figure axis must be checked before comparing magnitudes.

If $\bar e$ is the error mean and
$s_e^2=N^{-1}\sum_i(e_i-\bar e)^2$, direct expansion gives
\[
\mathrm{RMSE}^2=\bar e^2+s_e^2.
\]
RMSE combines systematic bias and scatter.
It is not the maximum error. For example, 99 exact predictions and one
1 nm error have RMSE 0.1 nm over 100 cases.

\section{Mean absolute percentage error}
The paper's Eq.~(3) is
\begin{equation}
\mathrm{MAPE}
=\frac{100\%}{N}\sum_{i=1}^N
\left|\frac{\hat y_i-y_i}{y_i}\right|.
\label{eq:MAPE}
\end{equation}
The absolute value makes each term nonnegative, including for negative
reference values. The denominator makes the statistic dimensionless.
It also gives large weight to small reference magnitudes and is undefined
when $y_i=0$.

For a 14 nm CD, an error of 0.014 nm is 0.1 percent.
For a 0.01 nm pattern shift, the same 0.014 nm error is 140 percent.
The second percentage looks alarming, but the physical displacement
error is identical. Conversely, a 1 percent error on a large quantity
can exceed a stringent absolute tolerance.

\figteach{percentage_error}{1.0}{A fixed absolute error of 0.01 nm
produces an arbitrarily large percentage as the reference shift approaches
zero. The right panel shows that the physical error does not diverge.}

\section{Why near-zero shifts are especially awkward}
Zero is often the desired shift. A metric that becomes undefined exactly
at the target is a poor sole measure of control quality.
The paper separates values around 0.14 nm in Fig.~8 and around 0.1 nm
in Fig.~10. The 0.14 nm scale is $1\%$ of the nominal 14 nm width;
it is not the paper's final $\pm0.4\nm$ acceptance boundary.

The text describes values above and below these thresholds but does not
fully document how signed shifts, exact zeros, or removed values were
handled. A careful replication would define whether the split uses
$|y_i|$, explicitly state any denominator floor, and report the sample
count in each group.

A worsening MAPE for smaller $|y|$ can occur with unchanged absolute
error. Therefore the graph alone does not demonstrate that a neural
network intrinsically cannot predict small numbers.
Noise, feature scaling, cancellation between coefficients, weak
observability, target density, and preprocessing can all contribute.
The paper's explanation should be treated as a suggested account,
not a mathematical inevitability of neural networks.

\section{Better companion measures for a placement budget}
Keep RMSE and add interpretable quantities:
\begin{align}
\mathrm{MAE}&=\frac1N\sum_i|e_i|,\\
e_{\max}&=\max_i|e_i|,\\
\mathrm{pass\ fraction}&=
\frac{\#\{i:|\mathrm{PS}_{y,i}^{\rm verified}|\leq b\}}{N}.
\end{align}
For a specified tolerance $b>0$, a useful normalized error is
$|e_i|/b$. Its denominator remains meaningful at zero target shift.
Quantiles show tail behavior; signed mean error detects systematic bias.
None of these removes the need to state which patterns and operating
conditions are included.

\worked{A candidate near the boundary}
Suppose desired shift is 0.395 nm and the prediction-to-simulation
error is $+0.020\nm$. The verified shift is 0.415 nm and fails a
0.4 nm limit. A small absolute discrepancy is consequential here.
The same discrepancy at a target of zero would still meet the limit.
The relevant question is error relative to remaining margin, not just
whether a number sounds small.

\section{Coefficient error versus image error}
For a held-out original pair $(\vct c,\vct y=F(\vct c))$,
the network outputs $\hat{\vct c}=G(\vct y)$.
Two separate residuals are
\[
\vct r_c=\hat{\vct c}-\vct c,\qquad
\vct r_y=F(\hat{\vct c})-\vct y.
\]
Figure 5 concerns coefficient differences. Figures 6 and 7 concern the
imaging differences after forward simulation.
A sizeable $\vct r_c$ can coexist with a small $\vct r_y$ if the two
wavefronts are nearly equivalent for the selected indicators.
However, coefficient disagreement by itself does not prove a valid
alternative: the forward residual must actually be checked.

Agreement in a few scalar metrics also does not establish identical
full contours, NILS, stochastic defect behavior, or performance on
unseen pattern classes. The paper validates the indicators it measures.

\section{What an independent simulation test establishes}
The 118 additional target vectors provide a useful test beyond the model's
training/validation set, as described by the authors.
Their proposed coefficients are sent through the same type of forward
simulator under the same lithographic conditions.
This tests consistency with the simulated model and conditions.
It is not an experimental measurement on a fabricated projector or wafer.

Using the same simulator for training labels and final checking can
validate the inverse approximation while leaving a common forward-model
bias invisible. Actual engineering deployment also needs calibrated
mask, resist, optical, and metrology models.
These limitations define the scope of the evidence rather than negate
the usefulness of a simulation-based study.

\part{Reading the results and constructing a defensible budget}
\chapter{Read the paper with this chapter beside it}\label{ch:walkthrough}
\section{Title, abstract, and the actual claim}
The title joins three ideas: an aberration budget, an EUV projector, and
the image quality of a contact layer. The budget is conditional on that
layer's selected patterns and imaging indicators. The paper does not
claim to characterize every possible mask pattern.

The abstract's phrase back propagating neutral network should be read as
\emph{backpropagation neural network}. The learned map takes imaging
indicators as inputs and proposes wavefront coefficients as outputs.
Forward re-simulation tests the resulting images.
The important conclusion is that the authors obtained candidate
wavefront distributions with useful simulated imaging performance under
their chosen conditions.

The phrase wavefront aberration margin should not immediately be read
as a certified independent tolerance on each mirror or each Zernike
coefficient. A margin needs a defined region, perturbation model, and
worst-case or probabilistic guarantee. The paper primarily demonstrates
successful coefficient combinations.

\section{Introduction: why this work was attempted}
\subsection{Shorter wavelength makes a fixed OPD more serious}
The 193 nm versus 13.5 nm comparison is worked through in
Chapter \ref{ch:waves}. In physical length units, the error is unchanged.
In waves and radians, it is about 14.3 times larger at EUV wavelength.
The consequences listed in the paper have distinct meanings:
\begin{itemize}
\item Reduced process window means fewer focus-dose combinations meet
all imaging tolerances.
\item Lower image fidelity means the printed contour departs more
from its intended shape.
\item Lower NILS means a shallower relative edge slope and generally
greater edge sensitivity.
\item Pattern centre offset is a placement error even if width is unchanged.
\item Unwanted artifacts can include distorted, bridged, or otherwise
undesired image features; the paper does not provide a separate
artifact-rate study.
\end{itemize}

\subsection{The earlier aberration-control methods}
The Introduction discusses several classes of prior work.
They are context, not additional algorithms implemented and validated
in this paper.

\textbf{Pattern-specific wavefront adjustment} uses available optical
controls to improve selected pattern images. The Nikon TAO example
illustrates that an optimum for one pattern need not minimize a
pattern-independent wavefront norm.

\textbf{Thermal prediction and control} models how absorbed power changes
optical shape or index and adjusts compensators before or during exposure.
A simple control description is
$\vct c(t)=\vct c_{\rm thermal}(t)+B\vct u(t)$,
where $\vct u$ represents actuator commands and $B$ their wavefront
influence. Predictive control uses a future estimate of
$\vct c_{\rm thermal}$ when choosing $\vct u$.
This equation is explanatory, not the controller in any cited instrument.

\textbf{Wavefront manipulators} provide practical degrees of freedom to
alter the pupil phase. The Introduction's FlexWave example concerns DUV
immersion equipment. It should not be read as a statement that arbitrary
25-dimensional phase changes are independently available in an EUV scanner.

\textbf{Optical prescription optimization} adjusts mirror shapes,
locations, and aspheric coefficients. Grouping methods divide a
complicated optical system into tractable subassemblies.
Ray tracing calculates actual surface intersections and directions.
Finite-aberration methods evaluate finite departures beyond a solely
infinitesimal paraxial analysis. These approaches operate on physical
design parameters rather than directly on an unconstrained image-to-mode map.

\textbf{Global search heuristics} mentioned in the references include
simulated annealing and particle swarm methods. Annealing sometimes
accepts a worsening objective by a temperature-dependent rule to escape
local minima. Particle swarm search updates candidate locations using
their own and the group's best positions.
The present paper's stated training algorithm is LM; it does not
report applying those design algorithms in its own 118-candidate test.

The quoted $\lambda/50$, $\lambda/70$, and $\lambda/30$ wavefront scales
belong to the cited contexts. At 13.5 nm they are 0.270, 0.19286, and
0.450 nm respectively. Comparisons require the same RMS definition,
removed modes, aperture weighting, field points, and wavelength.

\section{Figure 1: the optical system}
\paperloc{PDF p.~1 / journal p.~7270.}
Read the drawing as a map of functional stages:
illumination creates the source distribution at the mask; the reflective
mask diffracts the field; projection optics form the wafer image.
The colored boxes identify these components, not separate mathematical
wavefront-error budgets.

The EUV system is reflective, unlike an ordinary camera drawn with glass
lenses. The naming issue in the caption was noted in Chapter \ref{ch:euv}.
Nothing in the figure gives a mirror prescription, surface sensitivity
matrix, or actuator range. Consequently it cannot turn predicted
Zernike coefficients into a unique hardware design.

\section{Equation (1): the entire wavefront represented by coefficients}
\paperloc{PDF p.~2 / journal p.~7271.}
The equation
\[
W(\rho,\theta)=\sum_{i=1}^{\infty}Z_iF_i(\rho,\theta)
\]
is a basis expansion. Its technical content is explained in
Chapter \ref{ch:zernike}. A finite implementation replaces the infinite
sum by a selected set of terms. Truncation neglects all remaining
wavefront structure; it does not prove that those higher terms vanish
in a real projector.

The coefficients are not independent physical defects merely because
the basis is orthogonal. A single mirror alignment error may excite
several coefficients, and a single coefficient may receive contributions
from many surfaces.
The missing basis normalization and index dictionary are necessary for
quantitative comparisons with another optical program.

\section{Table 1 and Figure 2: illumination and process conditions}
\paperloc{PDF p.~3 / journal p.~7272.}
Table 1 fixes wavelength, numerical aperture, incidence geometry,
source shape, nominal stack, focus range, and coefficient sweep.
Together these define the conditional forward map.
Figure 2 shows the angular distribution of illumination in normalized
pupil coordinates. Bright regions are active source directions.
They are not four holes in the physical mask and not four printed contacts.

The source is separated from the projection pupil: the source chooses
incident directions, and the projection pupil selects diffracted outgoing
components. Chapter \ref{ch:partial} derives how they interact.
At the end of the paper the authors state that this source was not well
optimized for the contact layer. That limitation applies to the learned
aberration preferences as well.

\section{Table 2 and Figure 3: what the mask represents}
\paperloc{PDF p.~3 / journal p.~7272.}
The top view in Fig.~3 depicts the opening and absorber in a periodic
mask cell. The side view depicts absorber topography on a multilayer.
The square example is 14 nm wide on 38 nm pitch in wafer-referred units.
The diagram makes it clear that a hole in the absorber exposes a
reflective region; it is not a transparent hole through an ordinary lens mask.

The absorber and multilayer phases affect the complex diffraction orders.
Their geometry is one reason horizontal and vertical imaging can differ.
Table 2 gives two feature families and through-pitch examples.
The conflict with Table 3 must be remembered whenever translating
a plot's Pattern\_1 through Pattern\_7 into an actual mask layout.

\section{Figure 4: the simplified network}
\paperloc{PDF p.~3 / journal p.~7272.}
Each circle is a scalar neuron and each connection represents a trainable
weight. The arrows of evaluation run from imaging metrics to aberration
coefficients. During training, derivative information moves backward
through the computational graph.

The figure is explicitly simplified. The reported architecture in
Section 2C has seven hidden layers, while the schematic shows fewer
hidden groups. Use the written widths for the literal parameter-count
calculation, while recognizing the implementation details remain incomplete.
The diagram provides no activation functions, output constraints, or
training/test split.

\section{Equations (2) and (3): comparing errors}
\paperloc{PDF p.~4 / journal p.~7273.}
Equation (2) is RMSE; Eq.~(3) is MAPE. Their derivations and
limitations are in Chapter \ref{ch:errors}.
Every time an error is quoted, identify the target quantity, physical
units, aggregation set, and treatment of zeros.

An RMSE across coefficients is not an RMSE across printed CDs.
A percentage error in a small signed shift is not interchangeable
with a percentage error in a 14 nm width.
These distinctions explain why several later plots can look very different
without a logical contradiction.

\section{Table 3: labels for all subsequent bar charts}
\paperloc{PDF p.~4 / journal p.~7273.}
REC and SQ are described as rectangular and square patterns, with
seven variants each. The pitch assignments conflict with Table 2,
as transcribed in Chapter \ref{ch:data}.
Until this is resolved, a comparison such as REC versus SQ can be
reported using the authors' labels but should not support a precise
claim about a particular reconstructed pitch family.

The bars retain information by pattern and direction rather than averaging
all of them into one number. That preserves, for example, a poor result
for one sparse pitch that an overall average might conceal.

\section{Figure 5: different coefficients for similar requested images}
\paperloc{PDF p.~4 / journal p.~7273.}
Each bar summarizes disagreement between a predicted coefficient and
the original simulator coefficient associated with the input image metrics.
The vertical axis is in wavelengths. Larger bars mean the inverse model
has not simply recovered the original wavefront.

The authors interpret this as evidence of alternative aberration
combinations. That interpretation is plausible only together with the
subsequent image re-simulation. By itself, Fig.~5 is also compatible
with inverse prediction error.
Nearly equivalent images may come from weakly observable combinations
or nonlinear inverse ambiguity, as discussed in Chapter \ref{ch:inverse}.

The plotted coefficient labels run from $Z_2$ to $Z_{25}$.
They reinforce the 24-varied-coordinate interpretation, but do not
resolve the network's 25-output statement.
The figure is not a measurement of mirror surface error or a tolerance
specification.

\section{Figure 6: physical errors after forward re-simulation}
\paperloc{PDF p.~5 / journal p.~7274.}
\textbf{Panel (a), CD:} the simulator is run with predicted coefficients,
and the resulting widths are compared with the input target widths.
Bars are separated into $x/y$ and REC/SQ. The vertical unit is nm.
The graph shows noticeably different errors by pattern family and pitch,
so an overall accuracy number would hide structure.

\textbf{Panel (b), PS:} this is the same comparison for feature-centre
shift. Direction matters. Several vertical-shift bars for the SQ-labelled
family are visibly larger than the corresponding horizontal-shift bars.
The physically relevant result is their absolute size in nm and their
relation to a specified placement tolerance.

\textbf{Panel (c), PVB:} the comparison concerns process-contour
variation, not nominal feature width. A good CD prediction does not
imply a good PVB prediction because the latter depends on several
focus-dose conditions.

The figure gives aggregate RMSE values, not maximum per-case residuals.
The paper does not release numerical arrays for these bars; fine numerical
values should not be inferred from raster-image heights.

\section{Figure 7: percentages tell a different story}
\paperloc{PDF p.~6 / journal p.~7275.}
Panels (a) and (b) show MAPE for CD and PVB. Panel (c) shows horizontal
PS, and panel (d) vertical PS. Most plotted errors for the first three
categories are relatively small percentages, while vertical PS has
substantially larger percentages.

This is the expected danger when $\mathrm{PS}_y$ is close to zero:
the reference denominator becomes small. The statement that most errors
are below 1 percent does not describe every bar, especially not the
vertical-shift category.
RMSE and MAPE summarize different functions of the residuals, so
their relative bar ordering need not match.

\section{Figure 8: splitting around 0.14 nm}
\paperloc{PDF p.~6 / journal p.~7275.}
Panel (a) concerns the larger-shift subset and panel (b) the smaller-shift
subset, using the stated 0.14 nm threshold.
The smaller-reference subset has much worse percentage errors.
The graph supports that descriptive statement.

To decide whether the underlying physical prediction is worse, one also
needs absolute errors, counts, and uncertainty within each subset.
The mathematical denominator effect alone can generate the observed trend.
The threshold is a diagnostic choice motivated by 1 percent of 14 nm;
it is not a pass/fail criterion for lithographic placement.

\section{Section 3: the additional 118 cases}
\paperloc{PDF pp.~5--7 / journal pp.~7274--7276.}
The authors refer to a 2021 IRDS white paper and infer a requirement
that vertical contact placement be within 1 nm. The paper does not give
an explicit allocation derivation from that roadmap to a standalone
projector-wavefront PS allowance. Treat 1 nm as the authors' stated
motivation rather than a universal independently established per-mode limit.

They then use a tighter target range, $-0.4$ to $+0.4\nm$, for vertical
shift in 118 sets of imaging indicators. The remaining indicators are
kept within training-data ranges. The targets originate in other
simulations under the same stated conditions, not in arbitrary independent
coordinate draws. The model produces 118 coefficient vectors and the
forward simulator checks their imaging.

\section{Table 4: the remaining target ranges}
\paperloc{PDF p.~6 / journal p.~7275.}
\begin{table}[htbp]\centering\small
\caption{Paper Table 4 transcribed without silently correcting its repeated
last header. All entries are nm. The final column is likely intended to
represent vertical PVB, but that is an inference.}
\begin{tabular}{@{}llrrrrr@{}}
\toprule Family & Bound & $\mathrm{CD}_x$ & $\mathrm{CD}_y$
& $\mathrm{PS}_x$ & PVB column 1 & PVB column 2\\ \midrule
REC & Maximum &29.636&27.447&4.233&2.888&2.676\\
REC & Minimum &26.927&25.520&2.319&1.850&1.881\\
SQ & Maximum &13.488&12.663&4.190&7.238&8.755\\
SQ & Minimum &12.023&9.7623&2.083&3.079&3.791\\
\bottomrule
\end{tabular}
\end{table}
The publication labels both final columns PVB\_X.
Using one as PVB\_Y is a reasonable hypothesis because six metrics include
both directions, but it is not confirmed.

The REC horizontal widths in this table are also near 27--30 nm despite
the nominal horizontal CD of 14 nm elsewhere. Printed CD can differ
from nominal CD, but the discrepancy and the earlier table conflict
prevent us from deciding whether this is a measurement convention,
strong image bias, an axis swap, or a publication error.
Do not relabel the numbers to force agreement.

These are training-range conditions, not independently justified
production acceptance limits for all five indicators.
In particular, allowing horizontal shifts of several nm within the
training range does not prove they satisfy a complete overlay budget.

\section{Figure 9: candidate ranges and accuracy}
\paperloc{PDF p.~7 / journal p.~7276.}
\textbf{Panel (a):} the minimum and maximum predicted values are plotted
for each coefficient. The authors identify small ranges for $Z_3$,
$Z_{20}$, $Z_{21}$, and $Z_{22}$ and connect them with vertical PS.
These are observations within their selected candidate population.
They do not by themselves identify independent sensitivities or prove
that each coefficient must always be tightly bounded.

The coefficient axis is in \emph{waves}. Some visible extrema are
roughly $-0.065\lambda$ or $+0.05\lambda$, well beyond the stated
training sweep of $\pm0.020\lambda$.
Those are approximate visual readings, not digitized data.
They flag an important issue: the inverse output may extrapolate beyond
the sampled coefficient range, or the publication may contain an
unexplained scale/reporting difference.
The paper does not establish a hard output-bound enforcement rule.

The caption's reference to the first 25 sets is also unclear in relation
to the 118 proposed combinations and 24 plotted coefficient indices.
It should not be converted into an undocumented sample count.

\textbf{Panel (b):} absolute RMSE of verified imaging indicators versus
target indicators is shown by pattern.
\textbf{Panel (c):} MAPE is shown for the indicators other than vertical
PS. \textbf{Panel (d):} vertical-PS MAPE is shown separately and is
large for some patterns. The paper text refers to panels (b) and (c)
when discussing PS percentage behavior, but the dedicated PS MAPE
panel is (d); use the actual axis and legend to interpret it.

The panel (a) intervals are \emph{marginal ranges across candidates}.
Taking the largest allowed value from every bar and combining them
does not generally produce another valid candidate.
Correlations between coefficients may be responsible for successful
compensation. Chapter \ref{ch:budget} provides an explicit counterexample.

\section{Figure 10: a second small-value split}
\paperloc{PDF p.~8 / journal p.~7277.}
This figure repeats the larger-versus-smaller vertical-shift comparison
for the new application cases using 0.1 nm instead of 0.14 nm.
The smaller-reference group again has larger MAPE.
The two split thresholds belong to different analyses; there is no
single fixed 0.14 nm filtering rule governing all later results.

Since the new targets were deliberately selected near zero, strong
percentage penalties are unsurprising. To judge a 0.4 nm budget, inspect
the signed verified values and worst-case absolute errors as well.
The figure does not provide those complete arrays.

\section{Figures 11 and 12: two individual candidates}
\paperloc{PDF pp.~8--9 / journal pp.~7277--7278.}
Each figure has four parts:
\begin{enumerate}
\item \textbf{(a) Coefficient vector:} a particular proposed wavefront,
plotted in waves. The vertical scales differ between the figures.
Compare numerical axes rather than visual bar height.
\item \textbf{(b) CD:} paired bars compare the relevant target/prediction
and simulator-reference widths for both families and directions.
\item \textbf{(c) Horizontal PS and PVB:} the same candidate is evaluated
for additional imaging requirements and process robustness.
\item \textbf{(d) Vertical PS:} the small signed shifts are shown
separately so that discrepancies are visible near zero.
\end{enumerate}
The prose defines the suffixes \texttt{\_pred} and \texttt{\_ref}.
Operationally, however, the network outputs \emph{coefficients};
it does not output the CD bars directly in the described architecture.
The meaningful comparison is the requested/input image vector against
the forward-simulated image vector for the proposed coefficients.
This distinction prevents misreading these graphs as an independently
trained forward imaging network.

Figure 11 includes small shifts near or across zero.
Figure 12 includes negative vertical shifts.
These signed examples reinforce why an absolute-value tolerance must
be used and why zero handling matters in MAPE.
The authors report that all 118 proposed combinations satisfy the
intended 0.4 nm vertical-shift condition after simulation.
The two illustrated examples support the qualitative comparison, but
are not a substitute for the unreleased full validation table.

\section{Conclusion and the limits explicitly acknowledged}
The authors conclude that the inverse model can propose aberration
distributions satisfying selected imaging requirements and can assist
budget analysis. They also state that the chosen source is not well
optimized for the contact layer and that the patterns were not corrected
by OPC. They propose including improved illumination and corrected
patterns in future work.

OPC changes $\widetilde t$ or the rigorous mask-order amplitudes.
Source optimization changes $S$. Both change the forward map $F$.
A coefficient combination preferred before those changes is not
guaranteed to remain preferred afterward.

The data-availability statement says the underlying data are not public
and may be obtained from the authors on request.
No such request has been made for this guide. The absent data limit
exact numerical reproduction; they do not prevent deriving the physics
or critically understanding the published experiment.

\section{The claim you should be able to explain accurately}
Under a fixed simulated EUV contact-layer setup, the authors train an
inverse neural model from selected image metrics to Zernike coefficients,
then verify proposed combinations by forward simulation. They find
alternative wavefronts with useful agreement in the selected indicators
and report acceptable vertical shifts for their 118 application cases.
This supports pattern-dependent aberration selection within the studied
setup. It does not establish a unique inverse, independent per-mode
tolerances, a manufacturable mirror prescription, or general production
performance across arbitrary layers.

\chapter{Turning candidate wavefronts into an engineering budget}\label{ch:budget}
\section{State the acceptance condition before assigning tolerances}
An engineering budget begins with an explicit list of requirements.
For pattern $p$, examples are
\[
|\mathrm{PS}_{y,p}|\leq b_{y,p},\quad
|\mathrm{CD}_{x,p}-\mathrm{CD}_{x,p}^{\rm target}|
\leq b_{{\rm CD},p},\quad
\mathrm{PVB}_{x,p}\leq b_{{\rm PVB},p}.
\]
The conditions should identify the field point, process range,
measurement definition, and whether compensation such as refocus or
global image registration is allowed.

The paper emphasizes a vertical-shift target near $\pm0.4\nm$ and keeps
other requested indicators in training-data ranges. That is a particular
candidate-generation experiment. A production budget would require
explicit physical limits for all relevant indicators, rather than
training coverage alone.

\section{A local constraint is a slab in coefficient space}
Let $\vct c_0$ be a verified candidate.
For indicator $i$, write its residual from the target as $r_i$ and its
gradient row as $\vct g_i^T$. Locally,
\[
r_i(\vct c_0+\delta\vct c)
\simeq r_i+\vct g_i^T\delta\vct c.
\]
The two-sided acceptance condition is
\begin{equation}
|r_i+\vct g_i^T\delta\vct c|\leq b_i.
\label{eq:slab}
\end{equation}
This is the region between two parallel hyperplanes.
Intersecting all such slabs gives a local feasible polytope.
An optical or mechanical constraint on coefficients further restricts it.
This geometry makes clear why tolerances are coupled.

\section{The largest local RMS ball around a candidate}
Suppose an orthonormal coefficient norm represents the retained RMS.
For any $\norm{\delta\vct c}\leq R$, Cauchy--Schwarz gives
\[
|\vct g_i^T\delta\vct c|
\leq\norm{\vct g_i}R.
\]
A sufficient local guarantee is
$|r_i|+\norm{\vct g_i}R\leq b_i$ for every $i$.
Hence a radius satisfying
\begin{equation}
R\leq\min_i\frac{b_i-|r_i|}{\norm{\vct g_i}}
\label{eq:ballbudget}
\end{equation}
fits inside the linearized constraints, omitting rows with zero gradient
unless their residual already violates the limit.

This is a bound on \emph{perturbations around a chosen candidate}.
It differs from the candidate's own RMS and from a maximum admissible
absolute RMS anywhere in the feasible set.
A system can have an appreciable baseline wavefront and still have
very little remaining tolerance to perturbation.

\section{Independent box limits require a joint worst-case test}
Suppose each coefficient perturbation is independently allowed within
$|\delta c_j|\leq a_j$. Then
\begin{equation}
\max_{|\delta c_j|\leq a_j}
|\vct g_i^T\delta\vct c|
=\sum_j|g_{ij}|a_j.
\label{eq:boxsupport}
\end{equation}
The upper bound follows from the triangle inequality.
It is attained by choosing the sign of each $\delta c_j$ to align
with $g_{ij}$.
A safe linear box therefore requires
\begin{equation}
|r_i|+\sum_j|g_{ij}|a_j\leq b_i
\quad\hbox{for every required metric }i.
\label{eq:boxbudget}
\end{equation}
Checking one coefficient at a time while fixing the others cannot
establish this joint guarantee.

\worked{The error in treating marginal ranges as a budget}
Suppose a required shift is proportional to $c_1+c_2$ and must obey
$|c_1+c_2|\leq0.4$ in chosen units.
The two candidates $(1,-1)$ and $(-1,1)$ both give zero shift.
Their observed coordinate ranges are $[-1,1]$ for both coefficients.
Combining the two maximum values gives $(1,1)$, whose shift is 2
and fails badly.

\figteach{feasible_budget}{.78}{A correlated feasible region and two
feasible candidates. The dashed box contains all independent mixtures of
the candidates' coordinate ranges. Its upper-right corner is not feasible.
This is why the min/max bars in the paper's Fig.~9(a) are not independent
guaranteed tolerances. The additional diagonal bound in this illustration
only makes the plotted region finite.}

\section{Ellipsoidal budgets and covariance}
For an ellipsoidal perturbation region
$\delta\vct c^TQ\delta\vct c\leq1$ with positive definite $Q$,
set $\vct z=Q^{1/2}\delta\vct c$. Then
\[
\vct g_i^T\delta\vct c
=(Q^{-1/2}\vct g_i)^T\vct z.
\]
Maximizing over $\norm{\vct z}\leq1$ gives
\begin{equation}
\max|\vct g_i^T\delta\vct c|
=\sqrt{\vct g_i^TQ^{-1}\vct g_i}.
\label{eq:ellipsoid}
\end{equation}
This allows tolerances to be wider along insensitive combinations and
tighter along sensitive combinations.

For random perturbations with covariance $C_c$,
linear propagation gives
\begin{equation}
C_y=JC_cJ^T,\qquad
\sigma_{y_i}^2=\vct g_i^TC_c\vct g_i.
\label{eq:covprop}
\end{equation}
The covariance formula is a statistical statement, not a deterministic
guarantee. If Gaussian assumptions are justified, a multiple of
$\sigma_{y_i}$ can be related to a tail probability.
For non-Gaussian tails or systematic errors, that interpretation may fail.

\section{A mode-specific placement calculation}
For our unit-RMS vertical tilt,
$W=c_yZ_1^{-1}=2c_yv$.
Equation \eqref{eq:tilt} gives
\begin{equation}
|\mathrm{PS}_y|=\frac{2|c_y|}{\NA}.
\end{equation}
At $\NA=0.33$, a pure-tilt requirement
$|\mathrm{PS}_y|\leq0.4\nm$ implies
\[
|c_y|\leq0.066\nm
=4.88889\,\mathrm{m}\lambda
\quad(\lambda=13.5\nm).
\]
A 20 milliwave tilt would shift the image by about 1.63636 nm.
This is an exact translation result within the stated scalar
shift-invariant pupil convention. It is not a numerical bound on the
paper's $Z_3$ because its normalization and axis convention are not given,
and compensating modes or registration could change the relevant allowance.

\section{From wavefront coefficients to physical hardware}
Let $\vct q$ collect physical design or adjustment variables:
mirror decentres, tilts, surface modes, or actuator commands.
Linearizing the optical system gives
\begin{equation}
\delta\vct c=B\,\delta\vct q,\qquad
\delta\vct y\simeq JB\,\delta\vct q.
\label{eq:hardware}
\end{equation}
$B$ is a physical influence matrix. Its column space contains the
wavefront changes actually attainable to first order.
A candidate coefficient vector outside that attainable set cannot be
implemented by the available controls, even if its simulated image is good.

A practical adjustment problem might minimize
\[
\norm{Q_y^{1/2}[F(\vct c_0+B\delta\vct q)-\vct y_*]}^2
+\alpha\norm{L_q\delta\vct q}^2
\]
subject to displacement, surface-slope, stress, alignment, or actuator
bounds. The paper does not supply $B$, so it does not solve this
hardware-realization step.

\worked{An illustrative mirror-height allocation}
Assume six equal, statistically independent mirror contributions, all
with local incidence angle $6^\circ$, negligible coating contribution,
and total permitted RMS OPD $0.27\nm$.
Using Eq.~\eqref{eq:mirror},
\[
\sigma_{W,\rm total}^2
=6(2\sigma_h\cos6^\circ)^2.
\]
Thus
\[
\sigma_h=\frac{0.27}{2\cos6^\circ\sqrt6}
=0.05542\nm.
\]
This is about 55 pm per surface under the stated assumptions.
Six perfectly correlated equal wavefront contributions would add in
amplitude and give a different allocation,
$0.27/[12\cos6^\circ]\simeq0.02262\nm$.
Neither number is a recovered tolerance of the projector in the paper.
They demonstrate how geometry and correlation enter a surface budget.

\section{Nonlinear robustness and process variation}
The Jacobian approximations are local. If a bound $M_i$ on the
relevant Hessian norm is known throughout a perturbation region, Taylor's
theorem gives a remainder estimate
\[
|R_i|\leq\frac12M_i\norm{\delta\vct c}^2.
\]
A conservative ball test can then include
$|r_i|+\norm{\vct g_i}R+\tfrac12M_iR^2\leq b_i$.
Without such a bound, a local linear result needs nonlinear checking
over the proposed tolerance region.

Likewise, checking focus at $-15,0,+15\nm$ does not mathematically
prove every intermediate focus passes. The extremum of a nonlinear
CD or PS curve can lie between samples.
PVB envelopes involve maxima and minima, and can become nonsmooth
when the process condition defining an extreme changes.

Useful engineering checks include nominal cases, candidate corners,
random perturbations from an explicit distribution, and targeted searches
for the worst violation. A finite test is evidence about the tested
region and sampler; it is not automatically a continuous worst-case proof.

\section{Finite success counts and probability}
As a teaching calculation, suppose $N$ cases are independently drawn
from a fixed operational distribution and every case passes.
If the true failure probability is $p$, the probability of seeing no
failures is $(1-p)^N$.
Setting this to 0.05 gives the usual one-sided 95 percent upper limit
\[
p_{\rm upper}=1-0.05^{1/N}.
\]
For $N=118$, this is about 2.51 percent.
Even under those strong sampling assumptions, zero observed failures
would not prove zero failure probability.
The paper does not establish that its selected 118 targets are
independent random samples from a manufacturing distribution, so this
calculation is not a statistical confidence claim about its results.

\section{Optical determinism does not remove photon and resist noise}
The paper's reported metrics are deterministic simulation indicators.
EUV photon energy is $E_\gamma=hc/\lambda\simeq91.84$ eV at 13.5 nm.
For a deliberately illustrative dose of $20\,\mathrm{mJ/cm^2}
=200\,\mathrm{J/m^2}$ incident on a $14\nm\times14\nm$ area,
\[
N_\gamma\approx
\frac{200(14\times10^{-9})^2}{hc/(13.5\times10^{-9})}
\simeq2.66\times10^3.
\]
This is an incident photon estimate for uniform exposure, not the actual
absorbed count in a printed contact image.

For ideal Poisson counting, $\Var N=\E N$ and relative count noise is
$1/\sqrt N$, about 1.94 percent for this example.
Absorption, secondary-electron transport, resist chemistry, and
nonuniform intensity add further structure.
The IRDS discussion identifies stochastic defects and pattern control
as major lithographic concerns \cite{irds}.
A wavefront that passes deterministic CD and PS criteria does not
thereby establish an acceptable missing-contact probability.

\section{A defensible use of the paper's method}
Use the inverse network as a candidate generator.
Verify each candidate with the forward model and all required patterns.
Check coefficient-domain coverage, physical realizability, and robustness
to manufacturing and process uncertainty. Then express the final budget
as an explicit joint set, distribution, or perturbation allowance with
documented conditions.

This preserves the paper's useful idea: imaging performance can guide
wavefront choices more directly than a single RMS number.
It also supplies the additional steps needed to turn promising examples
into an engineering tolerance statement.

\chapter{A reproducible laboratory for the key ideas}\label{ch:lab}
\section{What the supplied calculation does}
The script \texttt{tools/reproduce\_examples.py} generates this book's
original figures and the files \texttt{numerical\_checks.txt} and
\texttt{teaching\_model\_results.json}.
It verifies disk-mode orthogonality and several analytic identities,
then calculates a simple periodic contact image.
It contains no trained version of the authors' network and does not
claim to reproduce the commercial reflective-mask simulation.

The dependencies are NumPy, SciPy, and Matplotlib.
The PDF and figures already included with this book can be read without
running Python.

\section{Declared optical and numerical assumptions}
The contact example uses:
\begin{itemize}
\item wavelength 13.5 nm, image NA 0.33, and an unobscured unit pupil;
\item a scalar binary thin mask containing one centred
$14\nm\times14\nm$ opening per $38\nm\times38\nm$ period;
\item four idealized source sectors, radii 0.6--0.8, centred at
$45^\circ,135^\circ,225^\circ,315^\circ$;
\item 60 equally weighted source samples: three radii uniformly spaced
in squared radius and five midpoint angular samples per sector;
\item exact analytic rectangle Fourier coefficients from
Eq.~\eqref{eq:rectcoeff}, retaining every order that can enter the pupil
for the selected source support;
\item pupil phase $2\pi W/\lambda$ and paraxial defocus phase
$-\pi z\NA^2\rho^2/\lambda$;
\item a fixed threshold chosen so that the no-aberration, zero-focus,
unit-dose reference has a 14 nm CD on its central cuts.
\end{itemize}
Equal spacing in squared radius approximates uniform annular area
weighting because $\rho\dd\rho=\tfrac12\dd(\rho^2)$.
These choices make the example reproducible; they are not inferred
hidden settings from the paper.

\section{How the calculation is organized}
For each source point:
\begin{enumerate}
\item Compute normalized order locations
$\vct u_{mn}=\vct u_s+(m/p_x,n/p_y)/f_c$.
\item Reject orders with $|\vct u_{mn}|>1$.
\item Multiply each retained mask coefficient by its aberration and
defocus phase.
\item Sum the complex plane-wave orders along an image cut.
\item Square the field magnitude.
\end{enumerate}
Average those intensities across source points.
Find left and right crossings of the fixed threshold by linear
interpolation on a dense spatial grid.
Then calculate CD and midpoint shift.
For pure tilt, orthogonal cuts pass through the analytically translated
feature centre; otherwise a fixed off-centre cut could falsely suggest
a width change even for a rigid translation.

\section{Same RMS, different images}
The three nonzero cases use one unit-RMS mode with
$c=0.27\nm=20\,\mathrm{m}\lambda$:
horizontal tilt, defocus, and cosine astigmatism.
All have the same piston-removed wavefront RMS in the declared disk basis.

\begin{table}[htbp]\centering\small
\caption{Computed teaching-model results, rounded for readability.
The exact arrays are saved with the source. These are not paper data.}
\begin{tabular}{@{}lrrrr@{}}
\toprule Case & CD$_x$ (nm) & PS$_x$ (nm)
& CD$_y$ (nm) & PS$_y$ (nm)\\ \midrule
Reference &14.00000&0&14.00000&0\\
Horizontal tilt &14.00000&$-1.63636$&14.00000&0\\
Defocus &13.94249&0&13.94249&0\\
Astigmatism &13.97132&0&13.97132&0\\
\bottomrule
\end{tabular}
\end{table}
Tilt agrees with the analytic displacement $-2c/\NA$ and preserves
the registered width. Defocus and astigmatism change the image profile
by different amounts.
In this particularly symmetric in-focus example, the astigmatism case
has equal central-cut CDs in the two axes. That is not a universal
statement about astigmatism; different focus, source, or mask conditions
can break this equality.

\figteach{contact_profiles}{1.0}{Original computed contact profiles.
Left: equal-RMS phase modes have different effects on a fixed-threshold
image. Right: a close-up of the right edge reveals the small change in
threshold crossing through focus. The $+15$ nm and $-15$ nm curves
coincide by symmetry in this example.
All curves use the same reference threshold rather than retuning dose
to conceal profile changes.}

\section{Focus and dose variation}
The demonstration samples focus at $-15,0,+15\nm$ and dose at
$0.95,1.00,1.05$. It collects left and right edge positions over all
nine combinations.
Both edge-band widths are approximately
\[
B_L=B_R=0.974586\nm
\]
in this symmetric teaching example.
The equality follows from the model symmetry, not from a universal
PVB identity. Because the centre remains fixed here,
the full CD range is twice one edge-band width.

\figteach{focus_dose_bands}{1.0}{Left: teaching Bossung curves computed
over a wider focus interval for illustration. Right: left and right
edge positions for the nine stated focus-dose samples. Shaded regions
show their individual process bands.}

\section{Independent checks built into the script}
The script uses Gauss--Legendre integration in $\rho^2$ and uniform
angular quadrature to check the Gram matrix of all 28 real modes
through radial degree six. Its largest numerical departure from
the identity is about $5.1\times10^{-14}$ in the supplied run.
It separately checks the balanced-quartic variance, the analytic
tilt displacement, preservation of registered widths under tilt, and
simple inverse-ambiguity and correlated-budget counterexamples.

It also prints the wavelength conversions, defocus RMS, 2033-row
arithmetic, hypothetical full-grid size, architecture parameter count,
and mask-scale consistency estimates. These checks test the calculations
actually used in this book.

The source sample count is intentionally modest for transparent,
quick execution. There is no claim that these pedagogical images are
electromagnetically rigorous or numerically converged to a production
accuracy requirement. Exact identities are checked tightly; approximate
optical-model output is labelled according to its assumptions.

\section{Experiments you can perform after reading}
\begin{enumerate}
\item Change pitch while keeping the hole width fixed. Track how the
admitted order set changes and identify which changes in the image
follow from altered interference.
\item Change the sign of a pure tilt coefficient. Confirm that the
shift reverses and registered widths stay fixed.
\item Add a small focus offset to astigmatism. Examine whether horizontal
and vertical CDs remain equal.
\item Reduce edge slope by changing source or threshold, then apply the
same dose perturbation. Compare measured edge motion with
Eq.~\eqref{eq:nilsdose}.
\item Increase source and spatial sampling. Check which plotted values
stabilize and how long the calculation takes.
\item Replace a single-mode perturbation by two modes together. Compare
the combined metric change with the sum of separate changes to reveal
nonlinear interactions.
\end{enumerate}
These exercises teach the mechanism behind the paper's dataset.
They do not require pretending that the missing commercial simulation
or neural-network weights are available.

\appendix
\chapter{Two derivations worth doing in full}\label{app:derivations}
\section{Deriving the Zernike radial formula and orthogonality}
Chapter \ref{ch:zernike} constructed low-order modes by subtracting
overlaps. Here is a general construction that proves the normalization
used throughout the book.

Take $m\geq0$, $k\geq0$, $n=m+2k$, and $t=\rho^2$.
Define
\begin{equation}
Q_k(t)=\frac{t^{-m}}{k!}
\frac{\dd^k}{\dd t^k}\left[t^{m+k}(t-1)^k\right],
\qquad R_{m+2k}^m(\rho)=t^{m/2}Q_k(t).
\label{eq:rodrigues}
\end{equation}
Although $t^{-m}$ appears, the derivative contains a factor $t^m$,
so $Q_k$ is an ordinary polynomial of degree $k$.
At $t=1$, only the term in which all $k$ derivatives act on
$(t-1)^k$ survives. Its value is $k!$, giving $Q_k(1)=1$
and $R_n^m(1)=1$.

Expand the bracket:
\[
t^{m+k}(t-1)^k
=\sum_{s=0}^k(-1)^{k-s}\binom ks t^{m+k+s}.
\]
Taking $k$ derivatives gives
\[
R_{m+2k}^m(\rho)
=\sum_{s=0}^k
\frac{(-1)^{k-s}(m+k+s)!}
{s!(k-s)!(m+s)!}\rho^{m+2s}.
\]
Set $j=k-s$, so that $m+2s=n-2j$ and $m+k+s=n-j$.
This becomes
\[
R_n^m(\rho)=\sum_{j=0}^{(n-m)/2}
\frac{(-1)^j(n-j)!}
{j![(n-m)/2-j]![(n+m)/2-j]!}\rho^{n-2j},
\]
which is Eq.~\eqref{eq:radial}.
The factorial formula is therefore the expanded form of a weighted
polynomial orthogonalization construction.

To prove orthogonality, take $\ell<k$ and evaluate
\begin{align}
\int_0^1 R_{m+2k}^mR_{m+2\ell}^m\,\rho\dd\rho
&=\frac12\int_0^1t^mQ_k(t)Q_\ell(t)\dd t\\
&=\frac1{2k!}\int_0^1Q_\ell(t)
\frac{\dd^k}{\dd t^k}[t^{m+k}(t-1)^k]\dd t.
\end{align}
Integrate by parts $k$ times. Boundary terms vanish: the bracket has
sufficient zeros at both $t=0$ and $t=1$.
The remaining integral contains $\dd^kQ_\ell/\dd t^k$, which is zero
because $Q_\ell$ has degree less than $k$.
This proves orthogonality for different radial degrees at fixed $m$.

For the norm, the leading coefficient of $Q_k$ is
\[
A_k=\frac{(m+2k)!}{k!(m+k)!}.
\]
Hence $\dd^kQ_k/\dd t^k=k!A_k$, and integration by parts gives
\begin{align}
\int_0^1[R_{m+2k}^m]^2\rho\dd\rho
&=\frac{A_k}{2}\int_0^1t^{m+k}(1-t)^k\dd t\\
&=\frac{A_k}{2}\frac{(m+k)!\,k!}{(m+2k+1)!}\\
&=\frac{1}{2(m+2k+1)}=\frac1{2(n+1)}.
\end{align}
The middle integral follows by repeated integration by parts:
$\int_0^1t^a(1-t)^b\dd t=a!b!/(a+b+1)!$ for nonnegative integer
$a,b$. Angular sine/cosine orthogonality supplies the remaining
normalization factors in Eq.~\eqref{eq:Zdef}.

\section{Why phase can move edges without changing total energy}
For $U=\FT^{-1}[\widetilde tP]$, Parseval's identity gives
\begin{equation}
\int|U(\vct x)|^2\dd^2x
=\int|\widetilde t(\vct f)|^2|P(\vct f)|^2\dd^2f.
\end{equation}
If only phase changes and $|P|$ stays fixed, the integrated image energy
is unchanged. For oblique illumination replace $P(\vct f)$ by
$P(\vct f+\vct s)$; the same conclusion holds for a fixed source point.
Source averaging preserves it.

Yet local intensity contains cross terms between frequencies.
Their phase differences change with $W$. Energy is redistributed,
altering maxima, minima, slopes, and threshold edges.
Phase aberration does not have to absorb light to degrade printing.

For periodic images, use energy per unit cell rather than an infinite
integral over repeated cells:
\[
\frac1{p_xp_y}\int_{\rm cell}|U|^2\dd x\dd y
=\sum_{m,n}|a_{mn}P_{mn}|^2.
\]
Distinct orders average to zero in the integrated cross terms,
but those cross terms remain essential locally.
If amplitude, aperture support, or coating reflectance changes too,
the total collected energy can change. The phase-only pupil model
isolates one mechanism rather than describing every physical perturbation.

\chapter{Exercises with worked solutions}\label{app:exercises}
Try each problem before reading its solution. Numerical examples use
this book's conventions unless explicitly stated otherwise.

\section{Phase, geometry, and diffraction}
\subsection*{1. Convert milliwaves into length and phase}
\textbf{Problem.} At 13.5 nm, convert $20\,\mathrm{m}\lambda$
into OPD and phase. Does this number alone specify RMS or PV?

\textbf{Solution.}
$W=13.5(20/1000)=0.270\nm$ and
$\phi=2\pi(0.020)=0.125664$ rad.
The conversion gives the size of the stated OPD quantity. Whether it
is a coefficient, RMS, PV, or local error must be specified separately.
For an RMS-normalized coefficient, its magnitude is its RMS contribution.

\subsection*{2. Convert a mirror-height error}
\textbf{Problem.} A local mirror height error is 0.10 nm at $6^\circ$
incidence from the normal. Estimate OPD and phase magnitude.

\textbf{Solution.}
$|\delta W|\simeq2(0.10)\cos6^\circ=0.19890\nm$.
At 13.5 nm this is about 14.734 milliwaves and 0.09257 rad.
The local geometric estimate excludes coating phase and changed ray
intersections.

\subsection*{3. Separate resolution and focus scaling}
\textbf{Problem.} At fixed wavelength, NA increases from 0.33 to 0.55.
How do $\lambda/\NA$ and $\lambda/\NA^2$ change?

\textbf{Solution.}
Their ratios are $0.33/0.55=0.60$ and
$(0.33/0.55)^2=0.36$.
The lateral diffraction scale decreases by 40 percent; the focus
scale decreases by 64 percent. Process-specific constants may also
change, so these are scaling comparisons.

\subsection*{4. Explain why one admitted order cannot draw a contact}
\textbf{Problem.} Show that one transmitted diffraction order has uniform
intensity. Why does the 38 nm example benefit from oblique light?

\textbf{Solution.}
With $U=a\ee^{\ii2\pi(f_xx+f_yy)}$, $|U|^2=|a|^2$ everywhere.
At the stated wavelength and NA, first-order normalized spacing for
38 nm pitch is 1.07656, outside the unit pupil for on-axis illumination.
Oblique source points shift several orders into the pupil, allowing
spatially varying interference cross terms.

\subsection*{5. Recover the phase-to-shift relation}
\textbf{Problem.} Two image orders separated by frequency $1/(38\nm)$
receive relative phase 0.10 rad. Find the fringe shift.

\textbf{Solution.}
$\Delta x=-p\Delta\phi/(2\pi)
=-38(0.10)/(2\pi)\simeq-0.60479\nm$.
This is a differential-phase result for two orders, not a general
coefficient-to-PS calibration.

\subsection*{6. Add amplitudes before intensities}
\textbf{Problem.} Two coherent equal-amplitude fields each have intensity
one. Find total intensity for relative phases 0, $\pi/2$, and $\pi$.
What if the sources are mutually incoherent?

\textbf{Solution.}
$I=2+2\cos\Delta\phi$, giving 4, 2, and 0.
For incoherent contributions the averaged cosine term vanishes and
intensity is 2. Averaging intensities prematurely loses interference.

\section{Wavefront statistics and Zernike modes}
\subsection*{7. Normalize defocus}
\textbf{Problem.} For $W=B\rho^2$, find its mean and piston-removed
RMS. Express the nonpiston part using $Z_2^0$.

\textbf{Solution.}
$\avg W=B/2$ and $\sigma_W^2=B^2(1/3-1/4)=B^2/12$.
The residual is
$B(\rho^2-1/2)=[B/(2\sqrt3)]Z_2^0$.
Its unit-RMS coefficient is $B/(2\sqrt3)$.

\subsection*{8. Distinguish raw and balanced spherical aberration}
\textbf{Problem.} Why is $B\rho^4$ not pure $Z_4^0$?
Find the best quadratic compensation in RMS.

\textbf{Solution.}
The raw quartic overlaps piston and defocus. For
$B\rho^4+d\rho^2$, minimizing piston-removed variance gives $d=-B$.
The residual after mean removal is
$B(\rho^4-\rho^2+1/6)=[B/(6\sqrt5)]Z_4^0$.
Its RMS is $|B|/\sqrt{180}$.

\subsection*{9. Count modes and total RMS}
\textbf{Problem.} How many real disk modes exist through degree six?
What RMS results if 24 orthonormal nonpiston coefficients are all
20 milliwaves?

\textbf{Solution.}
There are $\sum_{n=0}^6(n+1)=28$ modes including piston.
RMS is $\sqrt{24}(20)=97.9796$ milliwaves.
Neither calculation determines the paper's unspecified single-index ordering.

\subsection*{10. Find a pure-tilt placement allowance}
\textbf{Problem.} For $W=cZ_1^{-1}$ and NA 0.33, what $|c|$
corresponds to 0.4 nm vertical shift?

\textbf{Solution.}
$Z_1^{-1}=2v$ gives $|\Delta y|=2|c|/\NA$.
Thus $|c|=0.4(0.33)/2=0.066\nm=4.88889$ milliwaves at 13.5 nm.
This assumes the declared basis and no later removal of the translation.

\subsection*{11. Decide which error a Strehl calculation omits}
\textbf{Problem.} A report removes piston, tilt, and defocus before giving
a high Strehl estimate. Does it guarantee small absolute pattern shift
at a fixed wafer plane?

\textbf{Solution.}
No. Tilt changes absolute placement while preserving registered shape;
defocus compensation changes the observation plane.
The residual quality excludes those freedoms. Placement and fixed-plane
focus must be constrained separately.

\section{Image edges and process variation}
\subsection*{12. Distinguish CD change from PS change}
\textbf{Problem.} Left and right edges move by $-0.2$ and $+0.3$ nm.
Find width and centre changes.

\textbf{Solution.}
$\delta\mathrm{CD}=0.3-(-0.2)=0.5\nm$.
$\delta\mathrm{PS}=[0.3+(-0.2)]/2=0.05\nm$.
Most of the change is widening, with a small net translation.

\subsection*{13. Use the edge-slope formula}
\textbf{Problem.} At a bright feature's right edge, $I'=-0.1$ intensity
units per nm and $\delta I=+0.02$. Find edge displacement.

\textbf{Solution.}
$\delta x=-\delta I/I'=+0.2\nm$.
The edge moves outward. If the left slope is $+0.1$ with the same
perturbation, that edge moves $-0.2\nm$, so width increases 0.4 nm
and the centre stays fixed.

\subsection*{14. Check NILS with an exactly soluble image}
\textbf{Problem.} Let $I(x)=I_0\exp[-x^2/(2s^2)]$ and $T=I_0/2$.
Find nominal width, NILS, and width sensitivity to $\ln D$.

\textbf{Solution.}
The positive edge is $r=s\sqrt{2\ln2}$, width is $w=2r$, and
$|\partial\ln I/\partial x|=r/s^2$.
Thus $\mathrm{NILS}=2r^2/s^2=4\ln2$.
At dose $D$,
\[
w(D)=2s\sqrt{2\ln(2D)},\qquad
\frac{\dd w}{\dd\ln D}=\frac{2s}{\sqrt{2\ln(2D)}}.
\]
At $D=1$, this equals $2w/\mathrm{NILS}$, verifying the two-edge
formula.

\subsection*{15. Show that PVB is not always CD range}
\textbf{Problem.} A 14 nm feature translates between centres
$-0.3$ and $+0.3$ nm over process conditions. Find edge-band widths
and CD range.

\textbf{Solution.}
Each edge translates over 0.6 nm, so $B_L=B_R=0.6\nm$.
Every width remains 14 nm, so $B_{\rm CD}=0$.
Contour variation cannot be inferred from CD range alone.

\section{Data and neural-network reasoning}
\subsection*{16. Audit the sample count}
\textbf{Problem.} Why does the single-mode sweep have 984 rows but at
most 961 distinct wavefronts?

\textbf{Solution.}
$41(24)=984$ rows, but the all-zero vector appears 24 times.
Removing 23 extra copies leaves at most 961 distinct vectors.
Other groups can contain further duplicates.

\subsection*{17. Check normalization and extrapolation}
\textbf{Problem.} A feature trained over $[10,30]$ nm maps to $[-1,1]$.
Find normalized values at 20 and 35 nm.

\textbf{Solution.}
$\tilde x=2(x-10)/20-1$ gives 0 and 1.5.
The latter input is outside the observed range. Scaling does not clip
it or make extrapolation reliable.

\subsection*{18. Count a network's parameters}
\textbf{Problem.} Count parameters in a dense $84\to64\to25$ network
with biases.

\textbf{Solution.}
There are $(84+1)64=5440$ in the first layer and
$(64+1)25=1625$ in the second, totaling 7065.
The paper's longer literal architecture has 24,485.
Input entries are not themselves trainable parameters.

\subsection*{19. Calculate one LM step}
\textbf{Problem.} For $e(\theta)=2\theta-1$, start at $\theta=0$
with damping $\mu=1$. Find the LM step.

\textbf{Solution.}
$e=-1$, $J_e=2$, and
$(2^2+1)\Delta\theta=-2(-1)$.
Hence $\Delta\theta=0.4$ and the new residual is $-0.2$.
With $\mu=0$, the linear problem is solved exactly by $\Delta\theta=0.5$.

\subsection*{20. Explain a valid coefficient mismatch}
\textbf{Problem.} For $F(c_1,c_2)=c_1+c_2$, an original pair is
$(1,0)$ and the inverse proposes $(0,1)$. Is it necessarily wrong?

\textbf{Solution.}
Both produce output 1. Coefficient error is nonzero, but forward
consistency error is zero. An additional indicator such as $c_1-c_2$
would distinguish them. Equivalence depends on the selected requirements.

\subsection*{21. Explain an invalid mean of valid inverses}
\textbf{Problem.} For $F(c)=c^2$, targets $+1$ and $-1$ both yield
input $y=1$. What single prediction minimizes their average squared
coefficient loss, and is it feasible?

\textbf{Solution.}
Loss is $\hat c^2+1$, minimized at $\hat c=0$.
But $F(0)=0$, not 1. Inverse ambiguity can undermine regression even
when every training pair is physically valid.

\section{Validation and budgeting}
\subsection*{22. Compare percentage and absolute error}
\textbf{Problem.} Reference shift is 0.005 nm and prediction is 0.015 nm.
Find absolute and percentage error. Does the prediction meet a 0.4 nm
magnitude limit?

\textbf{Solution.}
Absolute error is 0.010 nm; percentage error is
$100(0.010/0.005)=200\%$.
The predicted 0.015 nm magnitude is below 0.4 nm.
A large percentage need not imply a physical limit violation.

\subsection*{23. Derive a safe two-coefficient box}
\textbf{Problem.} A local shift change is
$\delta y=2\delta c_1-\delta c_2$ and remaining allowance is 0.3.
Assign equal independent bounds $|\delta c_j|\leq a$.

\textbf{Solution.}
Maximum absolute change is $2a+a=3a$, so $a\leq0.1$.
Testing the first coefficient alone would allow 0.15; testing the
second alone would allow 0.3. Those individual limits cannot be
simultaneously permitted independently.

\subsection*{24. Propagate correlated uncertainty}
\textbf{Problem.} Use gradient $\vct g=(1,2)^T$ and covariance
$C=\begin{pmatrix}0.01&0.005\\0.005&0.04\end{pmatrix}\mathrm{nm}^2$
to find metric variance.

\textbf{Solution.}
$\vct g^TC\vct g=0.01+4(0.04)+4(0.005)=0.19\,\mathrm{nm}^2$.
Standard deviation is $0.43589\nm$.
Ignoring covariance gives $0.17\,\mathrm{nm}^2$ and underestimates
variance for this gradient.

\subsection*{25. Identify what remains unknown}
\textbf{Problem.} What essential information is missing for reproducing
every bar in the paper's Fig.~9 exactly?

\textbf{Solution.}
The coefficient convention; mask layouts resolving Tables 2 and 3;
source, stack, polarization, and simulator settings; metric extraction
definitions; training and test arrays; filtering and scaling settings;
network activations, weights, and split; and all 118 target and
verification vectors. The paper explains a workflow but does not
supply these artifacts. A reproduction must obtain or explicitly replace them.

\chapter{Reading map, notation, and source audit}\label{app:reading}
\section{Verified entry points into the supplied textbooks}
Printed page numbers refer to numbers on book pages.
PDF page numbers are one-based viewer positions in the supplied files.
The table gives verified entry points, not an assumed constant offset.

\begin{longtable}{@{}L{.27\textwidth}L{.44\textwidth}L{.20\textwidth}@{}}
\toprule Topic & Section and printed page & PDF page\\ \midrule\endhead
Scalar waves & Born--Wolf \S1.3, p.~14 &43\\
Reflection/refraction & Born--Wolf \S1.5, p.~36 &65\\
Layered media & Born--Wolf \S1.6, p.~51 &80\\
Characteristic matrices & Born--Wolf \S1.6.2, p.~55 &84\\
Periodic layers & Born--Wolf \S1.6.5, p.~66 &95\\
Geometrical-optics limit & Born--Wolf \S3.1, p.~109 &138\\
Stops and pupils & Born--Wolf \S4.8.2, p.~186 &215\\
Wave/ray aberrations & Born--Wolf \S5.1, p.~203 &232\\
Diffraction integral & Born--Wolf \S8.3.1, p.~375 &413\\
Circular-aperture diffraction & Born--Wolf \S8.5.2,
formula on p.~396 &434\\
Microscope resolution & Born--Wolf \S8.6.3, p.~419 &459\\
Aberration diffraction & Born--Wolf \S9.1.1, p.~460 &559\\
Reference-sphere changes & Born--Wolf \S9.1.2, p.~462 &557\\
Small-error intensity & Born--Wolf \S9.1.3, p.~463 &556\\
Zernike definition & Born--Wolf \S9.2.1, p.~464 &555\\
Radial-polynomial table & Born--Wolf \S9.2.1, p.~465 &554\\
Wavefront expansion & Born--Wolf \S9.2.2, p.~466 &553\\
Aberration tolerances & Born--Wolf \S9.3, p.~468 &551\\
Small-error criterion & Born--Wolf \S9.3, p.~469 &550\\
Coherent image transfer & Born--Wolf \S9.5.1, p.~483 &534\\
Incoherent transfer & Born--Wolf \S9.5.2, p.~486 &531\\
Mutual intensity & Born--Wolf \S10.4.1, p.~507 &510\\
Extended incoherent source & Born--Wolf \S10.4.2, p.~508 &509\\
Hopkins propagation & Born--Wolf \S10.4.2(b), p.~512 &505\\
\midrule
Wave-optics postulates & Saleh--Teich \S2.1, p.~40 &62\\
Complex wave amplitude & Saleh--Teich \S2.2, pp.~41--42 &63--64\\
Interference & Saleh--Teich \S2.5, p.~58 &80\\
Spatial harmonics & Saleh--Teich \S4.1, p.~105 &127\\
Angular spectrum & Saleh--Teich \S4.1, pp.~106--107 &128--129\\
Optical Fourier transform & Saleh--Teich \S4.2, p.~116 &138\\
Diffraction & Saleh--Teich \S4.3, p.~121 &143\\
Airy diffraction & Saleh--Teich \S4.3, p.~123 &145\\
Image formation & Saleh--Teich \S4.4, p.~127 &149\\
Coherent transfer function & Saleh--Teich \S4.4, p.~130 &152\\
Defocused pupil & Saleh--Teich \S4.4, p.~133 &155\\
Electromagnetic theory & Saleh--Teich \S5.1, p.~152 &174\\
Reflection/polarization & Saleh--Teich \S6.2, p.~209 &231\\
\bottomrule
\end{longtable}

For a first pass, use Saleh--Teich for waves and Fourier imaging, then
Born--Wolf for wave aberrations, Zernike theory, and partial coherence.
Born--Wolf uses older notation and electromagnetic unit conventions in
places. This guide uses SI for Maxwell's equations and explicitly
declares its complex-phase and polynomial normalization conventions.
Do not equate symbols across books by appearance alone.

The statistical-optics chapter in Saleh--Teich is absent from the
provided file. Its Fourier-transform appendices are also outside the
supplied excerpt. The needed definitions and derivations are included here.

\section{Symbol dictionary}
\begin{longtable}{@{}L{.22\textwidth}L{.53\textwidth}L{.16\textwidth}@{}}
\toprule Symbol & Meaning here & Units\\ \midrule\endhead
$\lambda,\lambda_0$ & Vacuum wavelength; also the wavelength in
vacuum image space &length\\
$k_0$ &$2\pi/\lambda_0$ &inverse length\\
$n,\widetilde n$ &Real or complex refractive index &dimensionless\\
$U,\vct U$ &Scalar or vector complex field amplitude &field scale\\
$I$ &Time-averaged intensity, often normalized &power/area or relative\\
$\mathcal L$ &Optical path length $\int n\dd\ell$ &length\\
$W$ &Actual-minus-reference phase-equivalent OPD &length\\
$\phi$ &Phase $2\pi W/\lambda$ &radians\\
$\sigma_W$ &Piston-removed RMS, unless other removals are stated &length\\
$\rho,\theta$ &Normalized pupil radius and azimuth &dimensionless; rad\\
$u,v$ &Normalized Cartesian pupil coordinates &dimensionless\\
$a$ &Physical pupil radius; opening width in rectangle examples
when stated locally &length\\
$\vct f$ &Transverse spatial-frequency vector &inverse length\\
$f_c$ &Image coherent cutoff $\NA/\lambda$ &inverse length\\
$P$ &Complex coherent pupil transfer function &relative amplitude\\
$S(\vct s)$ &Source weight density &reciprocal coordinate area\\
$Z_n^m$ &This guide's RMS-normalized real basis function &dimensionless\\
$c_n^m$ &Coefficient multiplying that function &length\\
$Z_i,F_i$ &Paper notation: coefficient and basis function &its convention\\
$z$ &Image-space focus displacement &length\\
$D$ &Relative dose multiplier in the contour example &dimensionless\\
$T$ &Fixed exposure threshold &exposure scale\\
$\vct y$ &Vector of imaging indicators &usually nm\\
$F,G$ &Forward imaging map and learned inverse &by arguments\\
$J$ &Imaging Jacobian $\partial\vct y/\partial\vct c$
&metric / coefficient\\
$J_e$ &Training-residual Jacobian versus network parameters &training scale\\
$\Theta$ &All trainable weights and biases &parameter-dependent\\
$\mu$ &LM damping; $\mu_{12}$ denotes coherence elsewhere &context-dependent\\
$B$ &Physical-variable-to-wavefront influence matrix &coefficient / control\\
\bottomrule
\end{longtable}

\section{Publication issues that must remain visible}
\begin{longtable}{@{}L{.31\textwidth}L{.63\textwidth}@{}}
\toprule Issue & Consequence\\ \midrule\endhead
Unspecified Zernike ordering and normalization
&Named high-index modes, RMS conversion, and transfer to another
simulator cannot be fixed uniquely.\\
$Z_2$--$Z_{25}$ versus 25 outputs
&There are 24 listed varied coordinates; the extra output is unexplained.\\
Tables 2 and 3 disagree
&Plot labels cannot be converted into a unique mask prescription.\\
Table 3 first SQ pitch pair
&$(56,56)$ also differs from Table 2's rectangular $(56,70)$.\\
Table 4 repeats PVB\_X
&Interpreting the second as PVB\_Y is plausible but unconfirmed.\\
Table 4 REC horizontal CDs
&Values near 27--30 nm differ strongly from the nominal 14 nm
horizontal size; the cause is unresolved.\\
CRAO and source shaping mixed in prose
&Chief-ray geometry and off-axis source distribution are distinct.\\
Figure 9 extrema exceed training sweep
&Output extrapolation or a scale issue needs clarification;
hard coefficient-box enforcement is not demonstrated.\\
Small-value filtering description
&Retained data, threshold units, and zero treatment are not reproducible.\\
Training split and activations absent
&Layer widths alone do not specify a trained model.\\
Metric extraction not fully defined
&Different CD/PS/PVB conventions can give different bars.\\
Roadmap-to-PS allocation not derived
&The 1 nm motivation is not a universal projector allowance.\\
Data and network weights absent
&The guide reproduces its teaching calculations, not the authors'
exact dataset.\\
\bottomrule
\end{longtable}

\section{A checklist for your own future experiment}
Record wavelength; image and object NA; magnification; physical and
wafer-referred geometry; pupil axes and support; OPD sign and units;
Zernike functions and normalization; removed modes and reference focus;
source weights; polarization; mask stack and materials; image or resist
model; contour measurement rules; dose-focus sampling; coefficient bounds;
training design and duplicate handling; split, scaling, and filtering;
network activations, loss, stopping, and random seed; and forward
verification of each candidate.

That record turns an interesting calculation into something another
researcher can examine, repeat, and improve.

\backmatter
\begin{thebibliography}{9}
\addcontentsline{toc}{chapter}{Bibliography and source notes}
\bibitem{chen}
Z. Chen, L. Dong, H. Ding, and Y. Wei,
\emph{Aberration budget analysis of EUV lithography from the imaging
performance of a contact layer in a 5 nm technology node},
Applied Optics \textbf{62}(27), 7270--7279 (2023).
\href{https://doi.org/10.1364/AO.496461}{doi:10.1364/AO.496461}.
Primary research source: the supplied 10-page
\texttt{Abberation Budget paper.pdf}. All three numbered equations,
four tables, and twelve figures were inspected.

\bibitem{bw}
M. Born and E. Wolf, \emph{Principles of Optics:
Electromagnetic Theory of Propagation, Interference and Diffraction of Light},
fourth edition, Pergamon Press.
The supplied 859-page scanned file was used, including the table of
contents, relevant physical-optics pages, and the reversed page run
containing Zernike and partial-coherence material.
Verified page locations are in Appendix \ref{app:reading}.
The fourth-edition preface is dated May 1968.

\bibitem{st}
B. E. A. Saleh and M. C. Teich, \emph{Fundamentals of Photonics},
second edition, Wiley (2007).
The supplied file is named \texttt{Saleh \& Teich Foundations of Photonics.pdf}
but the title identifies \emph{Fundamentals}. The excerpt ends at printed
page 278; references here are to content present in that excerpt.

\bibitem{asml}
ASML, \emph{TWINSCAN NXE:3400C}.
Official description used for the NXE system designation, reflective
4:1 reduction optics, and 0.33 NA.
\url{https://www.asml.com/en/products/euv-lithography-systems/twinscan-nxe3400c}.
Accessed 24 September 2026. This supplementary verification does not
replace the paper's simulation specification.

\bibitem{irds}
IEEE, \emph{International Roadmap for Devices and Systems: 2021 Update,
Lithography}.
\url{https://irds.ieee.org/images/files/pdf/2021/2021IRDS_Litho.pdf}.
Used for the distinction between node labels and dimensions and the
context of stochastic and pattern-control requirements.
The paper's reference 20 instead points to the companion LER white paper,
\url{https://irds.ieee.org/images/files/pdf/2021/LER_white_paper.pdf}.
The 1 nm vertical-PS inference is attributed to the paper's authors;
its exact allocation was not independently established from that companion.

\bibitem{mapminmax}
MathWorks, \emph{mapminmax: Process matrices by mapping row minimum
and maximum values to [-1 1]}.
\url{https://www.mathworks.com/help/deeplearning/ref/mapminmax.html}.
Official documentation for rowwise scaling and reuse of stored settings.
Accessed 24 September 2026.

\bibitem{trainlm}
MathWorks, \emph{trainlm: Levenberg--Marquardt backpropagation}.
\url{https://www.mathworks.com/help/deeplearning/ref/trainlm.html}.
Official documentation for damping, squared-error training, and
validation-failure stopping. Accessed 24 September 2026.
This guide explains the paper's historical method rather than
recommending a particular current MATLAB API.
\end{thebibliography}

\end{document}
